\chapter{Erklärung alterthümlicher und dunkler Ausdrücke}

\quellseite{36}

Aesern (nach Luther's Schreibart: euern), Spr. 17, 9.; wiederholen, wiedererregen.

Astersabath, Luc. 6, 1.; der erste Sabbath nach dem zweiten Tage des Osterfestes.

Aufmutzen, Sir. 13, 27.; hervorziehen, ans Licht bringen, groß machen.

Aufsetzen, 2 Kön. 18, 29.; 19, 10.; 2 Chron. 32, 15.; überreden, betrügen.

Aufsetzen, Mich. 12, 9.; aufstacheln.

Aufzug, 1 Macc. 8, 26.; Ausflucht.

Barte, Ps. 74, 6.; ein breiteres Beil, eine Art.

Crethi und Plethi, 2 Sam. 8, 18.; Läufer und Scharfrichter, die Leibwache Davids.

Dreiling, Jes. 40, 12.; ein drei Finger breites Maß.

Ehrlich, 4 Mos. 16, 2.; Luc. 14, 8.; vornehm, angesehen.

Endlich, Spr. 21, 5.; 22, 29.; Luc. 1, 39.; was zu Ende kommt und zu Ende bringt, eilig, emsig, tüchtig.

Enthalten, Ps. 51, 14.; aufrecht halten. Sich enthalten, Apok. 1, 13.; sich aufhalten.

Erwegen (erwägen), sich des Lebens, Weish. 17, 15.; 2 Cor. 1, 8.; das Leben verloren gehen.

Fast, sehr, überaus.

Ferge, Hes. 27, 27.; Fährmann, Schiffer.

Feser (Fäser), Jes. 5, 7.; die Haare oder Fasern an den keimenden Reben oder andern Pflanzen, welche in die Erde gesenkt, zuerst Wurzel fassen.

Fladdernholz, Hes. 27, 5.; Mastholz.

Geilen, Luc. 11, 8.; auf eine unverschämte Art, mit geifernder Stimme, um etwas betteln.

Geren, Hes. 16, 8.; Hagg. 2, 13.; die Schleppe des Kleides, der Schooß; auch der Flügel, der Zipfel des Mantels.

Glum, Jes. 32, 2.; schlammig, unrein; im Gegensatz von lauter, B. 14.

Gnätze, 3 Mos. 14, 56.; Krätze.

Greten, Hes. 16, 25.; aufspreizen, die Füße auseinander sperren.

Hellig, Jer. 2, 25.; athemloß, keuchend, lechzend.

Hengel, 1 Kön. 22, 34. Nach Luther's Erklärung: „Da das Schwert anhängt, von der Achsel über bis auf die Hüfte.“ Andere erklären: die Gelenke und Fugen am Panzer.

Hümpler, Spr. 26, 10.; ein Stümper, schlechter Arbeiter.

Igel (Egel), Spr. 30, 15.; der Blutigel.

Keulig, 1 Kön. 7, 41.; kugelig, rund.

Köden, Jes. 28, 7.; herausreißen.

Kogel, Hes. 23, 15.; langherabhängende Kopf binden.

Kolk, 3 Mos. 11, 36.; eine Grube, worin sich Wasser sammelt.

Koller, Jes. 3, 23.; Halskragen, Brusttuch; auch Leibgürtel, Apok. 19, 12.

Kollern, 1 Sam. 21, 13.; sich unsinnig und toll gebärden, sich auf dem Boden herumwälzen.

Koller, 2 Kön. 8, 15.; Decke, Bettdecke.

Krebs, Weish. 5, 19.; Eph. 6, 14.; 1 Thess. 5, 8.; Brustharnisch, Panzer (von der krebsförmigen Gestalt so genannt).

Kreuel, 2 Mos. 27, 3.; 1 Sam. 2, 13.; Fleischgabel.

Legel, 1 Sam. 16, 20.; Jer. 13, 12.; 48, 12.; nach Adelung: ein rundes, hölzernes Gefäß, in Gestalt einer Tonne.

Leibzucht, 1 Macc. 10, 54.; Leibrente, Unterhalt auf Lebenszeit.

Librarei, 2 Macc. 2, 13.; Bücherei, Büchersammlung.

Läden (läden, nach Luther's Schreibweise: leden), Hiob 21, 11.; Ps. 29, 6. u. ö.; springen, hüpfen; Apok. 9, 5.; mit den Füßen ausschlagen.

Löden, Hof. 7, 14.; heulen, plärren.

Lotterbube, Apok. 17, 18.; Schwätzer, unnützer Mensch.

Maul, 2 Kön. 13, 29.; 18, 9.; 1 Kön. 1, 33.; Maulthier.

Melechet des Himmels, Jer. 7, 18. u. ö.; Himmelskönigin, wie Luther in den ersten Ausgaben selbst übersetzt hatte.

Meye, Bar. 6, 8., im schlimmen Sinne; Richt. 5, 30., im guten, eine ehrliche Frauenperson.

Müller, Pred. 12, 3. 4.; Mühlmagd; die Stimme der Müllerin: der schwache Ton der Mühle.

Nehmen, 4 Mos. 11, 5.; eine Art Kürbis. (Nach dem Hebr.: Wassermelonen.)

Nadeln, 1 Kön. 7, 31.; Nadeln, Hesek. 1, 11.; Büchsen, erhabene Zierrathen von Metall, bes. auf Schilden.

Racha, Matth. 5, 22.; loser, nichtswürdiger Mensch.

Rachus, Jer. 15, 13.; 17, 3.; Hes. 23, 46.; Gegenstand feindlicher Plünderung.

Rechtfertigen, Apok. 12, 19.; vor Gericht laden, anklagen und bestrafen.

Saff, 2 Mos. 12, 8.; 4 Mos. 9, 11.; eine Brühe von bittern Kräutern.

Sangen, Ruth 2, 14.; 1 Sam. 17, 17.; gesengte, geröstete Aehren.

Schauer (Scheuer), Hiob 27, 18.; Obdach gegen böse Witterung.

Schemen, Ps. 39, 7.; Spr. 27, 19.; Schattenbild.

Scheuer, Sir. 50, 10.; ein schimmerndes Gefäß; Becher.

Schier, Ps. 119, 87.; Apok. 17, 15.; 1 Tim. 3, 14.; bald, schnell.

Schröter, Jer. 48, 12.; Leute, welche einen Stoff ausschroten, d. i. zertheilen und stückweise veräußern.

Seettig, 1 Tim. 6, 4.; (von Seuche) krank, siech.

Sichermahl, 1 Sam. 20, 20.; sicheres festes Ziel (im Gegensatz des Schießens ins Ungewisse).

Sinflut, 1 Mos. 7, 6.; Sir. 39, 27. u. ö.; große, weitverbreitete Flut.

Sponde, Amos 3, 12.; Bettgestelle.

Spunden, 1 Kön. 6, 9. u. ö.; decken, täfeln.

Stöcken, Weish. 2, 19.; in den Stock legen, zwingen (Ps. 105, 18.); peinigen.

Sträter, Hof. 6, 9. nach Luther's Erklärung: Strauchdiebe, Strauchräuber.

Tartische, 1 Kön. 10, 17.; ein langer Schild.

Tendchen, 5 Mos. 14, 5.; eine Art von Gemsen oder wilden Ziegen.

Thelding, a) Jagdübung, der Termin zur gerichtlichen Handhabung; b) die Verhandlungssache; dann unnützes Hin- und Herreden, Geschwätz. Hiob 35, 16.; Hes. 22, 28. Daher Narrenheiding Eph. 5, 4., loses, unnütziges Gerede.

Thramen, 1 Kön. 6, 6.; Balken.

Thürstig, thürstlich, 1 Mos. 34, 25.; Hiob 12, 6. u. ö.; kühn, beherzt, dreist.

Treusen, Mich. 2, 6.; schelten, strafen.

Unrath, Matth. 26, 8.; nach Luther's Erklärung: Es ist Schade drum, sie sind unräthlich damit umgegangen.

Verschleifen, 2 Sam. 14, 14.; schleifen, sich verlaufen.

Verschleißen, Sir. 14, 18.; zerreißen, verschleißen, verderben.

Versprechen, a) Richt. 9, 23.; verwünschen, verfluchen; b) Marc. 7, 2.; tadeln, sich über etwas aufhalten.

Wind, Spr. 30, 31.; Windhund, Windspiel.

Zauen, sich, 2 Sam. 5, 24.; sich sputen, eilen.

Zerlecken, zerlecken, Pred. 12, 6.; leck werden, Ritze bekommen.
% Scan 37: no retained apparatus prose; source blocks omitted.
% Scan 38: no retained apparatus prose; source blocks omitted.
% Scan 39: no retained apparatus prose; source blocks omitted.
\part{Das erste Buch Mose}

\chapter{Das 1. Kapitel}

\section{Vorrede}

\quellseite{40}

Aller guten und nützlichen Werke Anfang, Mittel und Ende stehet allein in Gottes Händen. Derowegen weil wir heute die Lection oder das Vorlesen der heiligen Biblia anfangen, so laßt uns herzlich beten, der allmächtige und gütige Gott wolle unsere Augen öffnen, und unsere Herzen erleuchten, daß wir die Wunder sehen in seinem Gesetz. Wie der königliche Prophet David im 119. Psalm redet; das ist: daß wir also sein heilig Wort täglich anhören, lernen, behalten und gebrauchen, auf daß wir seinen väterlichen Willen und seine Gnade und Barmherzigkeit daraus erkennen, und daß Glaube, Liebe, Trost, Geduld, Zucht, Gehorsam und Unterthänigkeit in uns allen dadurch gepflanzet werde, wachse und zunehme. Wer nun solches begehret, der spreche von Herzen, Amen. Das hilf uns allen, o lieber Gott und Vater, Amen. Darauf höret an in Gottes Namen das erste Kapitel des ersten Buchs Mose.

\section{Schöpfung der Welt}

a Geist) Wind ist zuweilen noch nicht gewesen, dann muß es den heiligen Geist deuten.

b Gut) das ist, nützlich, fein, löblich.
\quellseite{41}

c Unterthan) was ihr bauet und arbeitet auf dem Lande, das soll nur eigen sein, und die Erde soll euch hierin dienen, tragen und geben.

d Speise für die Menschen und Thiere etc.

\section{Summarium}

Ist eine Historia, wie Gott durch sein ewiges, allmächtiges Wort Himmel und Erden, sammt allen Kreaturen erschaffen, und am sechsten Tage den Menschen, der solches alles genießen und gebrauchen sollte, gemacht, und zum Herrn darüber gemacht habe. V. 27.

Auf daß wir lernen Gottes Güte in solchem Geschöpf erkennen, und unser Herz in allerlei Nöthen damit trösten, daß wir so einen allmächtigen und gnädigen Gott haben, der uns so reichlich versorget, und so väterlich mit allem, was wir bedürfen, bedacht hat. Solchen gnädigen Willen Gottes haben auch die Heiden aus der Schöpfung und Kreatur erkennet, und derselben Gott für gut und gnädig gehalten. Aber Gottes Wort hält den Christen einen höhern und größern Trost für, dabei sie müssen schließen, daß Gott gnädig sei; sintemal er nicht allein alles den Menschen zu Nutz erschaffen, sondern auch seinen eingebornen Sohn zum Opfer für die Sünde, und in den Tod hat gegeben. Darum sollte niemand an solches gnädigen Gottes Hülfe, Gunst und Rettung zweifeln, es sei gleich die Noth so groß sie immer wolle.

Merke hie den Artikel unsers Glaubens. Gott ist, ehe einige Kreatur ist, das ist, er ist ewig, und hat bei sich das Wort, das ist auch ewig und Gott gleich; denn es ist, ehe alle Kreaturen sind. Und weil durch solches Wort alle Kreaturen erschaffen werden, muß folgen, daß solches Wort allmächtig ist. Zum dritten nennt Moses hie den Geist des Herrn, der alles Geschöpf lebendig macht und erhält. Daß also Moses klar anzeiget, GOTT Vater, GOTT Sohn, oder das Wort, und GOTT der Heilige Geist sei der einige, ewige, allmächtige GOTT, der Anfang und das Ende, wie unser Christen-Glaube ausweiset.

\section{Gebet}

Wir danken Gott herzlich für die Erschaffung aller Kreaturen; der wolle ferner die Werke seiner Hände segnen und erhalten, und uns in seiner Gnade zum ewigen Leben bewahren, Amen.
\chapter{Das 2. Kapitel}

\section{Vorrede}

\quellseite{42}

Gott hat uns Menschen erschaffen, daß wir ihm sollen gleich sein in Gerechtigkeit, Heiligkeit, Barmherzigkeit, Wahrheit und in andern Tugenden, die in ihm wesentlich leuchten, und von welchen die heiligen zehn Gebote lehren. Aber, ach weh uns allen, wie ungleich sind wir Gott! Das lässet uns erkennen, und Gott abbitten, im wahren Glauben und Vertrauen auf das Versöhn-Opfer Jesu Christi, unsers HErrn. Wir wollen auch, so lange wir leben, unsere fürnehmste Sorge sein lassen, daß wir durch Hülfe des heiligen Geistes Gott, unserm Schöpfer, von Tage zu Tage gleichförmiger werden, und also gegen uns selbst züchtig, gegen unsern Nächsten gerecht, und gegen Gott gottselig leben, Tit. 3. Amen. Das soll auch erinnern das 2. Kapitel.

\section[Vom Sabbath, des Menschen wesentlichen Stücken, Paradiesgarten, Gottes]{Vom Sabbath, des Menschen wesentlichen Stücken, Paradiesgarten, Gottes Gebot, Ehestand}





a Pison) ist das große Wasser in Indien, das man Ganges heißt.

b Hevila) ist Indien-Land.

c Gihon) ist das Wasser in Aegypten, das man Nil heißt.


d Hiddekel) ist das Wasser in Assyrien, das man Tigris heißt.

e Euphrath) aber ist das nächste Wasser in Syrien, das man Euphrat heißt.




*Um ihn sey) das ist, kein Thier nahm sich des Menschen an, um ihn zu sein, daß ihm hülfe sich mehren und nähren etc.
\section{Summarium}

\quellseite{43}

Ist eine Historia, wie Gott durch sein ewiges, allmächtiges Wort Himmel und Erden, sammt allen Kreaturen erschaffen, und am sechsten Tage den Menschen, der solches alles genießen und gebrauchen sollte, gemacht, und zum Herrn darüber gemacht habe. V. 27.

Auf daß wir lernen Gottes Güte in solchem Geschöpf erkennen, und unser Herz in allerlei Nöthen damit trösten, daß wir so einen allmächtigen und gnädigen Gott haben, der uns so reichlich versorget, und so väterlich mit allem, was wir bedürfen, bedacht hat. Solchen gnädigen Willen Gottes haben auch die Heiden aus der Schöpfung und Kreatur erkannt, und derhalben Gott für gut und gnädig gehalten. Aber Gottes Wort hält den Christen einen höhern und größern Trost für, dabei sie müssen schließen, daß Gott gnädig sei; sintemal er nicht allein alles den Menschen zu Nutz erschaffen, sondern auch seinen eingebornen Sohn zum Opfer für die Sünde, und in den Tod hat gegeben. Darum sollte niemand an solches gnädigen Gottes Hülfe, Gunst und Rettung zweifeln, es sei gleich die Noth so groß sie immer wolle.

Merke hie den Artikel unsers Glaubens. Gott ist, ehe einige Kreatur ist, das ist, er ist ewig, und hat bei sich das Wort, das ist auch ewig und Gott gleich; denn es ist, ehe all Kreaturen sind. Und weil durch solches Wort alle Kreaturen erschaffen werden, muß folgen, daß solches Wort allmächtig ist. Zum dritten nennt Moses hie den Geist des Herrn, der alles Geschöpf lebendig macht und erhält. Das also Moses klar anzeiget, GOTT Vater, GOTT Sohn, der das Wort, und GOTT der heilige Geist sei der einige, ewige, allmächtige GOTT, der Anfang und das Ende, wie unser Christen-Glaube ausweiset.

\section{Beten}

Wir danken Gott herzlich für die Erschaffung aller Kreaturen; der wolle ferner die Werke seiner Hände segnen und erhalten, und uns in seiner Gnade zum ewigen Leben bewahren, Amen.

d Speise für die Menschen und Thiere etc.

c Unterthan) was ihr bauet und arbeitet auf dem Lande, das soll euer eigen seyn, und die Erde soll euch hierin dienen, tragen und geben.
\quellseite{44}
allen Kreaturen; da kann die Vernunft anders nicht richten, denn es sei keine Hülfe mehr im Himmel und auf Erden. Aber hier aus diesem Exempel sollen wir lernen, daß Gott auch helfen wird, wenn er gleich von allen Kreaturen verlassen sind. Doch thut er solche Hülfe um seines Sohnes willen, den er alhier dem menschlichen Geschlecht zu senden zugesaget etc.

d Derselbe hieß ist das erste Evangelium und Verheißung von Christo, geschehen auf Erden, daß er sollte Sünde, Tod und Hölle überwinden, und uns von der Schlangen Gewalt selig machen. Daran glaubet Adam mit allen seinen Nachkommen; davon er zum Christen, und selig worden ist von seinem Fall.

e (Stechen) plagen, kreuzigen und martern. Denn so gehets auch: Christus zertritt dem Teufel seinen Kopf, das ist, sein Reich, des Todes, Sünd' und Hölle. So sticht ihn der Teufel in die Fersen, das ist, er tödtet und martert ihn und die Seinen leiblich.

f Heva) Kai heißt Leben; daher kömmt Heva oder Hava, Leben oder Lebendige.

\section{Summarium}

Ist eine Historia, wie der Mensch durch den Teufel von Gottes Wort abgeführet, und in die Sünde und Strafe Gottes durch den Ungehorsam gefallen sei, uns zur Warnung, daß wir lernen fest an Gottes Wort halten, und uns vor Sünden hüten, welche vom Teufel ihren Anfang hat, und den Tod, sammt allem andern Jammer, über die Menschen gebracht hat. Dagegen soll man auch den Trost merken, welchen hie Gott Adam und Eva gibt, daß des Weibes Same, Christus, der Schlange den Kopf zertreten, das ist, dem Teufel seine Macht nehmen, und durch sein Leiden uns aus dem Tode und Sünde zur Gerechtigkeit und ewigen Leben helfen soll. Solchen Trost sollen wir fest behalten und uns das Fersenstechen des Teufels, das ist, allerlei Widerwärtigkeit dieses Lebens nicht ärgern lassen; denn im Paradies hat es angefangen, und muß für und für also gehen. Der Teufel gönnet uns die Seligkeit nicht; wiederum will Gott uns in Sünden nicht verderben lassen, und schaffet uns Hülfe durch seinen Sohn, der für uns am Kreuze stirbt und bezahlet. Solchen Trost und Hülfe sollen wir mehr achten, denn des Teufels Fluch und Feindschaft.

\section{Gebet}

Lob, Ehr, Preis und Dank sei Jesu Christo, unserm Heilande, der des Teufels Werk zerstöret hat; der erhalte uns im wahren Glauben und guten Gewissen zum ewigen Leben, Amen.
\chapter{Das 4. Capitel}

\section{Vorrede}

\quellseite{45}

Wie Kain ein Mörder worden; also ist kein Mensch auf Erden, er trägt bei sich ein Stück von einem Mörder. Denn gleich wie der Teufel ein Lügner und Mörder ist von Anfang, Joh. 8., das ist, von seinem Falle her: also hat er unter andern auch diese zween giftige Pfeile in aller Menschen Herzen, Fleisch und Blut geschaffen, und hat ihnen das Schandmal angehängt, daß sie nicht allein von Natur zu Lügen, sondern auch zu Zorn, Haß und Todtschlag geneigt sind. Derwegen ein jedes gute Achtung auf sich gebe, und lasse die Sünde ja nicht herrschen in seinem sterblichen Leibe, ihr Gehorsam zu leisten in ihren Lüsten. Röm. 6. Dies lehret auch das 4. Capitel.

\section{Kains Brudermord. Seine Nachkommen}

a Ei, Gott sei gelobt, da habe ich erlanget den HErrn, den Mann, den Samen, der dem Satan, der höllischen Schlange, den Kopf zertreten soll; der wird thun.

b Hebräisch lautet Thür so viel, als man das offen stehet oder aufgethan wird. Marc. 7, 34.: Sephahtah, thue dich auf etc. Und ist die Meynung: die Sünde lieget und ruhet, wie ein Docht lieget und ruhet. Aber sie lieget in der Thür, das ist, sie wird offen stehen, oder offenbar werden, ob der Sünder wohl eine Zeit lang ihrer dahin gehet, als sei die Sünde todt und schlaf.

c Redet mit Habel das ist, Scham halber mußte er sich äußerlich stellen und reden mit seinem Bruder, weil er gestraft ward, ob er wohl im Herzen ihn zu tödten gedachte. Also ist Kain aller Heuchler und falscher Heiligen Vater.
\quellseite{46}
*c. 5, 6.

d fing man an) nicht daß zuvor nicht auch Gottes Name wäre gepredigt; sondern nachdem durch Kains Bosheit der Gottesdienst gefallen war, ward er diesmal wieder aufgerichtet und irgend ein Altärlein gebauet, dahin sie sich versammleten, Gottes Wort zu hören und zu beten.

\section{Summarium}

Das ist eine Historia, wie der leibliche Segen mit Adam und Eva gegangen, und darneben die Sünde sich auch gewaltig an Kain erzeuget habe, der seinen Bruder Abel erschlug.

Solches Exempel gehet noch immerdar in der Welt. Denn der Abel, das ist die rechte christliche Kirche, die auf Gottes Wort und seine Barmherzigkeit in Christo sich gründet, muß sich von Kain, das ist, von der falschen Kirche und den Werktheiligen leiden. Aber es soll doch nicht ungerochen bleiben. Denn Gott will das unschuldige Blut suchen, und dräuet gar ernstlich, daß alle, so sich hier vergreifen, mit bösem Gewissen und allerlei Unglück ihr Leben lang sollen geplaget werden.

Hier ist sonderlich zu merken die Predigt, die (der HErr) Adam dem Kain thut, und warnet, er soll der Sünde den Raum nicht zu lang lassen, sondern an sich halten, sie würde ihn sonst verführen. Bist du fromm, spricht er, so bist du angenehm, das ist, halt dich in Gottesfurcht, und hüte dich für der Sünde, so wird Gott ob dir halten. Bist du aber nicht fromm, das ist, wenn du der Sünde folgest, und willst auf Gottes Furcht und Wort nicht sehen, so ruhet die Sünde, das ist, du hast Lust und Freude daran; aber es währet nicht lange, sie schläft in der Thür, an einem sehr unruhigen Ort, da muß sie aufwachen. Alsdann folget unermeßliche Angst, Schrecken, Jammer und der Tod selbst; wie man hier an Kain und an Judas stehet. Darum warnet (der HErr) Adam, laß du ihr nicht ihren Willen, das ist, wehre dich, widerstehe der Sünde. Und wie es Paulus heißt, würge sie. Da soll uns zu helfen des Weibes Saame, der will bei uns halten, und uns wider die Sünde helfen. Wir für uns würden zu schwach sein.

\section[Notum]{\unsicher{Notum}}

Helffe der treue Gott, daß ein jedes unter uns seine Natur und böse Zuneigung erkenne, die Sünde meide und meide, und dagegen christlich und ehrbar lebe, zu Lob und Preis seinem heiligen Namen, Amen.

\chapter{Das 5. Kapitel}

\section{Vorrede}

Wenn ein Mensch noch so lange lebet, so muß er doch endlich sterben. Dies lehren auch die gewöhnlichen Sprüchworte: Der Tod ist eine unvermeidliche Noth; Sterben ist unser Erbe; und: Wider Todes Gift und Macht, hilft keines Krautes Kraft und Saft. Derwegen laßt uns wachen und beten, damit, wenn wir auch durch den zeitlichen Tod aus dieser Welt scheiden, wir im Glauben und mit gutem Gewissen einschlafen, und am jüngsten Tage zum ewigen freudenreichen Leben mit allen Auserwählten von Todten auferstehen. Dazu höret an das 5. Kapitel.

\section[Geschlechtsregister der Patriarchen von Adam an bis auf]{Geschlechtsregister der Patriarchen von Adam an bis auf Noah}
\chapter{Das 4. Kapitel}

\section{Vorrede}

\quellseite{47}

Wie Kain ein Mörder worden: also ist kein Mensch auf Erden, er trägt bei sich ein Stück von einem Mörder. Denn gleich wie der Teufel ein Lügner und Mörder ist von Anfang, Joh. 8., das ist, von seinem Falle her: also hat er unter andern auch diese zween giftige Pfeile in aller Menschen Herzen, Fleisch und Blut geschaffen, und hat ihnen das Schandmaal angehängt, daß sie nicht allein von Natur zu Lügen, sondern auch zu Zorn, Haß und Todtschlag geneigt sind. Derwegen ein jedes gute Achtung auf sich gebe, und lasse die Sünde ja nicht herrschen in seinem sterblichen Leibe, ihr Gehorsam zu leisten in ihren Lüsten. Röm. 6. Dies lehret auch das 4. Kapitel.

\section{Kains Brudermord. Seine Nachkommen}

a Ei, Gott sei gedankt, da habe ich erlanget den HErrn, den Mann, den Samen, der dem Satan, der höllischen Schlange, den Kopf zertreten soll; der wirds thun.

b Hebräisch lautet Thür so viel, als was offen stehet oder aufgethan wird, Marc. 7. 34.: Hephatha, thue dich auf etc. Und ist die Meinung: die Sünde lieget und ruhet, wie ein Ochs lieget und ruhet. Aber sie lieget in der Thür, das ist, sie wird offen stehen, oder offenbar werden, ob der Sünder wohl eine Zeit lang sicher dahin gehet, als sei die Sünde todt und schlafe.

* Redet mit seinem Bruder) das ist, Scham halber mußte er sich äußerlich stellen und reden mit seinem Bruder, weil er gestraft ward, ob er wohl im Herzen ihn zu tödten gedachte. Also ist Kain aller Heuchler und falscher Heiligen Vater.
\quellseite{48}

d Fing man an) nicht daß zuvor nicht auch Gottes Name wäre gepredigt; sondern nachdem durch Kains Bosheit der Gottesdienst gefallen war, ward er dazumal wieder aufgerichtet und irgend ein Altärlein gebauet, dahin sie sich versammleten, Gottes Wort zu hören und zu beten.

\section{Summarium}

Das ist eine Historia, wie der leibliche Segen mit Adam und Eva gegangen, und darneben die Sünde sich auch gewaltig an Kain erzeuget habe, der seinen Bruder Abel erschlug.

Solches Exempel gehet noch immerfort in der Welt. Denn der Abel, das ist die rechte christliche Kirche, die auf Gottes Wort und seine Barmherzigkeit in Christo sich gründet, muß sich von Kain, das ist, von der falschen Kirche und den \unsicher{Werth[e}iligen] leiden. Aber es soll doch nicht ungerochen bleiben. Denn Gott will das unschuldige Blut suchen, und dräuet gar ernstlich, daß alle, so sich hier vergreifen, mit bösem Gewissen und allerlei Unglück ihr Leben lang sollen geplaget werden.

Hier ist sonderlich zu merken die Predigt, die (der HErr) Adam dem Kain thut, und warnet, er soll die Sünde den Zaum nicht zu lang lassen, sondern an sich halten, sie würde ihn sonst verführen. Bist du fromm, spricht er, so bist du angenehm, das ist, halt dich in Gottesfurcht, und hüte dich für der Sünde, so wird Gott ob dir halten. Bist du aber nicht fromm, das ist, wenn du der Sünde folgest, und willst auf Gottes Furcht und Wort nicht sehen, so ruhet die Sünde, das ist, du hast Lust und Freude daran; aber es währet nicht lange, sie schläft in der Thür, an einem sehr unruhigen Ort, da muß sie aufwachen. Aldann folget unermäßliche Angst, Schrecken, Jammer und der Tod selbst; wie man hier an Kain und an Judas siehet. Darum warnet (der HErr) Adam, laß du ihr nicht ihren Willen, das ist, wehre dich, widerstehe der Sünde. Und wie es Paulus heißt, würge sie. Da soll uns zu helfen des Weibes Same, der will bei uns halten, und uns wider die Sünde helfen. Wir für uns würden zu schwach sein.

\section{Beten}

Helfe der treue Gott, daß ein jedes unter uns seine Natur und böse Zuneigung erkenne, die Sünde meide und meide, und dagegen christlich und ehrbar lebe, zu Lob und Preis seinem heiligen Namen, Amen.

\chapter{Das 5. Capitel}

\section{Vorrede}

Wenn ein Mensch noch so lange lebet, so muß er doch endlich sterben. Dies lehren auch die gewöhnlichen Sprichworte: Der Tod ist eine unvermeidliche Noth; Sterben ist unser Erbe; und: Wider Todes Gift und Macht, hilft keines Krautes Kraft und Saft. Derwegen laßt uns wachen und beten, damit, wenn wir auch durch den zeitlichen Tod aus dieser Welt scheiden, wir im Glauben und mit gutem Gewissen einschlafen, und am jüngsten Tage zum ewigen freudenreichen Leben mit allen Auserwählten von Todten auferstehen. Dazu höret an das 5. Capitel.

\section[Geschlechtsregister der Patriarchen von Adam an bis auf]{Geschlechtsregister der Patriarchen von Adam an bis auf Noah}
\quellseite{49}
wegen laßt uns in diesen gefährlichen Zeiten fleißig beten, und ein jedes lebe in seinem Beruf und Stande, in wahrer Buße christlich und unsträflich. Wird ja Gott in folgender Zeit die Gottlosen mit den wohlverdienten Strafen heimsuchen, so wird er auch die Gläubigen und Frommen in solchem allgemeinen Uebergange in Gnaden schützen, trösten und zum ewigen Leben bewahren. Davon höret auch an das 7. Kapitel.

\section{Die Sündflut bricht ein}








\section{Summarium}

Ist eine Historia, wie Gott um der Sünde willen durch die Sündflut alles vertilget, was sich auf Erden regt, B. 12.-17. An dem frommen Noah aber beweiset er seine Gnade und erzeiget sich so gnädig und väterlich gegen ihn, daß er ihn mit den Seinen nicht allein erhält im Kasten, sondern auch für sie sorget, daß sie zu essen haben, und die Creaturen, so dem Menschen zu seinem Brauch dienstlich waren, erhalten werden.

Wie nun dies eine Sündflut heißt, darum, daß die Sünde also gräulich gestraft ist: also heißt die heilige Taufe eine rechte Sündflut, durch
\quellseite{50}
welche die Sünde nicht wie dort gestraft, sondern erlauft und vergeben wird durch den Glauben an Christum, welcher die geistliche Arca ist, daran wir erhalten werden.

\section{Psalm}

Der starke barmherzige GOtt tröste und errette alle, die in Noth sind, lindere uns auch die wohlverdienten Strafen, und schütze und erhalte uns gnädiglich in diesen gefährlichen Zeiten, Amen.

\chapter{Das 8. Kapitel}

\section{Vorrede}

Ach wie oft hätte uns Gott auch strafen müssen, wenn er nach unserm Verdienst mit uns hätte handeln sollen! Wer erkennet aber diese Gottes Gnade, Güte und Barmherzigkeit? Dergemahnet Gott, daß er sich so gnädig und väterlich gegen uns erzeigt, und daß er die wohlverdienten Strafen aufschiebt. Lasset uns auch in wahrer Buße leben. Denn auf muthwillige Sünden folget doch endlich große Strafe. Denn Gott kann bald also zornig werden, als gnädig er ist, Sirach 5. So hat kein Zorn über die Gottlosen ein Aufhören; also erinnern wir uns heute abermals der Sündflut. Davon auch lehret das 8. Kapitel.

\section{Die Sintflut nimmt ein Ende}

a Oelblatt) das Blatt bedeutet das Evangelium, das der heilige Geist in der Christenheit hat predigen lassen; denn Oel bedeutet Barmherzigkeit, und Friede, davon das Evangelium lehret.
\quellseite{51}
wegen laßt uns in diesen gefährlichen Zeiten fleißig beten, und ein jedes lebe in seinem Beruf und Stande, in wahrer Buße christlich und unsträflich. Wird ja Gott in folgender Zeit die Gottlosen mit den wohlverdienten Strafen heimsuchen, so wird er auch die Gläubigen und Frommen in solchem allgemeinen Uebergange in Gnaden schützen, trösten und zum ewigen Leben bewahren. Davon höret auch an das 7. Kapitel.

\section{Die Sintflut bricht ein}

\section{Summarium}

Ist eine Historia, wie Gott um der Sünde willen durch die Sintflut alles vertilget, was sich auf Erden regt, B. 12.–17. In dem frommen Noah aber beweiset er seine Gnade und erzeiget sich so gnädig und väterlich gegen ihn, daß er ihn mit den Seinen nicht allein erhält im Kasten, sondern auch für sie sorget, daß sie zu essen haben, und die Kreaturen, so dem Menschen zu seinem Brauch dienstlich waren, erhalten werden.

Wie nun dieß eine Sintflut heißt, darum, daß die Sünde also gräulich gestraft ist: also heißt die heilige Taufe eine rechte Sündflut, durch
\quellseite{52}
welche die Sünde nicht wie dort gestraft, sondern erlöst und vergeben wird durch den Glauben an Christum, welcher die geistliche Arca ist, daran wir erhalten werden.

\section{Psalm}

Der starke barmherzige GOtt tröste und errette alle, die in Noth sind, lindere uns auch die wohlverdienten Strafen, und schütze und erhalte uns gnädiglich in diesen gefährlichen Zeiten, Amen.

\chapter{Das 8. Capitel}

\section{Vorrede}

Ach wie oft hätte uns Gott auch strafen müssen, wenn er nach unserm Verdienst mit uns hätte handeln sollen! Wer erkennet aber diese Gottes Gnade, Güte und Barmherzigkeit? Derwegen danket Gott, daß er sich so gnädig und väterlich gegen uns erzeigt, und daß er die wohlverdienten Strafen aufschiebt. Lasset uns auch in wahrer Buße leben. Denn auf muthwillige Sünden folget doch endlich große Strafe. Denn Gott kann bald also zornig werden, als gnädig er ist, Sirach 5. So hat sein Zorn über die Gottlosen kein Aufhören; also erinnern wir uns heute abermals der Sündflut. Davon auch lehret das 8. Capitel.

\section{Die Sündflut nimmt ein Ende}



a Oelblatt) das Blatt bedeutet das Evangelium, das der heilige Geist in der Christenheit hat predigen lassen; kein Oel bedeutet Barmherzigkeit, und Friede, davon das Evangelium lehret.
\section{Das 1. Buch Mose. 8. 9}

\quellseite{53}

\section{Summarium}

Ist eine schöne Historia, wie die Sintflut wieder abgenommen und Noah sammt den andern, die mit ihm waren, und allen Thieren frisch und gesund aus dem Kasten gingen, und Gott den Segen auf ein neues beide über Menschen und Vieh ausgesprochen habe; durch welchen Segen alles, bis an der Welt Ende, erhalten würde. Sonderlich aber ist dies Exempel wohl zu merken, daß Noah dem HErrn opfert und für seine Wohlthat danket. Und Gott solch Dankopfer ihm läßt so wohl gefallen, daß er verheißet, er wolle fortan seinen Zorn nimmer so heftig über alles Fleisch gehen lassen, sintemal der Mensch von Natur böse und durch die Sünde also verderbet ist, daß er seiner Natur halben nicht anders kann, denn Unrecht thun; solches beweget den gnädigen Gott nicht zum Zorn noch Ungnade, sondern daß er sich unser erbarmet und Sünde vergibt.

\section{Votum}

GOtt der Vater aller Barmherzigkeit vergebe uns alle unsere Sünde, und erhalte uns in aller Gefahr und Noth, mit Leib und Seele zum ewigen Leben, um JEsu Christi seines lieben Sohnes willen, Amen.

\chapter{Das 9. Kapitel}

\section{Vorrede}

Die heiligen Engel im Himmel und alle Creaturen und Geschöpfe am Firmament, in der Luft und auf Erden, bleiben in ihrer Ordnung, zu der sie sind erschaffen, und sind Gott gehorsam: allein der Mensch, der doch die edelste Creatur ist, ist GOtt ungehorsam; wie solches GOtt selber bezeuget im 14. Psalm, da er also saget: Sie sind alle abgewichen und allesamt untüchtig, da ist keiner, der Gutes thue, auch nicht einer. Derwegen läßt uns lernen, was da sei des HErrn Wille, auf daß wir also leben nach seinem Wort durch Hülfe des heiligen Geistes, damit wir auch in Gottes Gnade und Segen leben, zeitlich und ewig, um unsers HErrn Jesu Christi Verdienst und Fürbitte willen. Dazu höret ferner an das 9. Capitel.

\section{GOtt bestätiget seinen Bund durch den Regenbogen}

a Durch Menschen) hier ist das weltliche Schwert eingesetzt, daß man die Mörder tödten soll.
\section{Summarium}

\quellseite{54}

Dies ist eine Historia, wie Gott insonderheit wiederholet den Segen über den Menschen, und setzet ihn zum Herrn über alle andere Kreatur, daß er derselben mächtig sein und zu seinem Nutzen brauchen soll.

Allein verbietet er das Blut, zum Zeichen, daß man Menschenblut nicht vergießen soll; denn hier stehet das Recht, welches Gott selbst spricht: Wer Menschenblut vergießt, deß Blut soll wieder durch Menschen, das ist, durch die Obrigkeit vergossen werden.

Darnach bekräftiget Gott durch das Zeichen des Regenbogens seine Verheißung, daß er fortan dergleichen nicht brauchen wolle.

\section{Der andere Theil des 9. Capitels}

\section{Summarium}

Dies ist ein Exempel, aus welchem Kinder und Unterthanen lernen sollen, wie sie ihre Eltern und Obrigkeit in allen Ehren sollen haben; auf daß der Fluch nicht auf sie erbe, welchen Noah über seinen Sohn Ham hie spricht, darum, daß er ihn unehret. Denn es ist nicht allein Sünde, daß er seine Lust an dem hat, daß er den Noah also entblößet siehet, der böse Bube will auch die andern Brüder dahin bringen, daß sie des alten Vaters lachen, ihn verspotten und nichts auf ihn halten noch geben sollen. Denn Ham thut zweierlei Sünde: den Vater verachtet er und hält sich für weiser und klüger. Da ist unmöglich, daß in einem solchen Herzen könnte einiger rechter Gehorsam sein, wie das Gott von allen Kindern fordert, und greuliche Strafe wider die ergehen läßt, so ihn nicht leisten.

Dagegen aber, wer seine Eltern und Oberherrn in Ehren hat, soll sich des Segens trösten, welchen Noah über seine anderen zween Söhne, Sem und Japheth, spricht. Denn Gott will es mit geistlichem und leiblichem Segen bezahlen, wie das vierte Gebot auch verheißt, daß Gott solchen Gehorsam und Ehrerbietung mit langem glückseligen Leben bezahlen wolle.

\section{Votum}

Helfe Gott uns allen, daß wir gegen unsere Obrigkeit und Eltern uns also verhalten, und in Gottesfurcht also unser Leben beschließen, auf daß wir für Christi, unseres HErrn, Richterstuhl wohl bestehen, Amen.

\chapter{Das 10. Capitel}

\section{Vorrede}

Dieses Kapitel wird nicht gelesen, denn es ist nur ein Geschlechts-Register und Namen-Buch der Nachkommen des Noah, welches dazu dienet, daß man daraus sehe die Fortpflanzung des frommen und gottesfürchtigen Geschlechts Noahs und der Verheißung seines Samens. Wiewohl, was von Nimrod gesagt wird, zu dem Ende zu merken ist, weil er, nachdem er das Reich und Gewalt überkommen, zu einem Tyrannen worden und die rechtgläubige Kirche sehr geplaget hat. Denn daß er hat gehabt den Namen eines gewaltigen Jägers für dem HErrn, ist nichts anderes, denn daß er die rechtgläubige Kirche unterdrückt habe, und nur darauf sei bedacht gewesen, wie er seine Macht und Gewalt vermehren und erweitern möge.
\section{Geschlecht und Nachkommen der Söhne Noah}

\quellseite{55}

a) Peleg auf deutsch eine Zertheilung.

\chapter{Das 11. Kapitel}

\section{Vorrede}

Vom heiligen alten Tobia stehet geschrieben Tob. 14.: Je älter er worden, daß er je mehr in Gottesfurcht habe zugenommen. Das Widerspiel stehet man heut am meisten Theil der Menschen. Desselbige aber ist dieses: Je älter die Leute werden, je ärger sie auch werden. Derwegen weil wir Gottes Wort haben, auch das täglich hören, so wollen wir auch in Kraft des heiligen Geistes nach Gottes Wort leben, wollen das Böse hassen und lassen, und das Gute lieben und üben, auf daß wir Gottes Zorn und Strafe auf uns nicht laden, sondern in seiner Gnade und unter seinem Schutze leben und sterben, Amen. Darauf höret nun an das 11. Kapitel.

\section{Thurm zu Babel. Verwirrung der Sprachen. Geschlecht Sems}
\quellseite{56}



a Babel) auf deutsch eine Vermischung oder Verwirrung.







\section{Summarium}

Dies ist eine Historia, wie die Welt, sobald sie ein wenig gewachsen, wieder sei böse, frech und hoffärtig worden, darum sie Gott den Ländern nach mit der Sprache getheilet und ihr Fürnehmen gehindert hat. Solche Theilung der Sprachen ist eine greuliche Sündenstrafe gewesen. Denn da ist so bald mit der Sprache das Predigtamt gefallen und ist so bald Abgötterei eingerissen. Denn ein jeder ihm einen eigenen Gottesdienst gemacht hat. In des einigen Sems Linie und Geschlecht ist die rechte Sprache, das Predigtamt und Gottesdienst blieben. Zum selben hat sich Japhet gehalten. Aber Ham, als der große Haufe, ist mit Macht in die Abgötterei gerathen. Gleich aber, wie Gott hier die Sünde strafet durch viele Zungen: also beweiset er am Pfingstage seine Gnade durch viele Zungen, daß man in allen Sprachen Gottes Gnade hört und das Predigtamt in allen Zungen gehet.

Darnach erzählet Mose das Geschlecht Sem, auf daß wir wissen die Ankunft des Erzvaters
\quellseite{57}
Abraham, da nachmals die Verheißung auf uns beruhet. Da stehet man, wie die Welt so bald abgenommen und die heiligen Leute nicht mehr so lange gelebet haben, wie Noah und seine Voreltern.

\section{Votum}

Behüte und regiere uns der treue GOtt alle, damit wir an Tugend und Ehrbarkeit nicht abnehmen, sondern zunehmen und zeitlich und ewig leben in Heiligkeit und Seligkeit, um Jesum Christi, unsers HErrn und Heilandes willen, Amen.

\chapter{Das 12. Kapitel}

\section{Vorrede}

Gleich wie die Sonne mitten im Winter und in der größten Kälte am weitesten von uns scheinet zu sein, da sie doch in ihrem Zirkel näher bei uns ist, als wenn die Hitze im Sommer am größten ist: also ist unser HErr und Helfer, der treue liebe Gott, uns auch am nächsten, wenn wir meinen, er sei gar weit von uns. Derowegen sagt man im gemeinen Sprüchwort: Je größer Noth, je näher Gott. Denn wenn Rath und Hilfe bei den Menschen aus ist, da lässet sich der große Nothhelfer und Rathgeber erst recht sehen mit Trost, Hilfe und Errettung. Davon lehret der König David im 34. Psalm also: Der HErr, spricht er, ist nahe bei denen, die zerbrochenes Herzens sind, er hilft denen, die zerschlagenen Gemüth haben. Des erfreuet und tröstet euch, ihr betrübte und nothleidende Herzen. Und höret darauf an das 12. Kapitel.

\section{Abram wird berufen und ziehet nach Egypten}
\quellseite{58}

\section{Summarium}

Hier gehet an die Legende des heiligen Patriarchen Abrahams, welchen Gott fordert aus dem Lande Haran ins Land Canaan, und verheißet ihm, er wolle ihn zum besondern Volk machen, und das Land Canaan eingeben; und zu seiner Zeit ihm einen solchen Samen erwecken, der ein gebenedeyter Same sei, das ist, durch welchen alle Welt aus dem Fluch zum Segen und Gottes Gnade soll gebracht werden.

Hie soll man erstlich lernen, daß Gott gnädig sei. Darnach soll man Acht haben auf das Exempel Abrahams, wie er Gott trauet und seinem Befehle folget und Gottes Wort prediget, und dennoch das Kreuz nicht außen bleibet, sintmal er Theurung halben in Egypten fliehen und da seines Weibes halben in Gefahr kommen muß. Aber weil er Gott glaubet und vertrauet, so hilft er ihm aus aller Noth; wie durch seine ganze Legende zu sehen ist; und strafet, die ihn beleidigen, wie er hie Pharao und sein Haus plaget, darum, daß er Sara ihm zum Weibe wollte nehmen.

\section{Votum}

Der gerechte starke Gott strafe noch heut alle muthwillige und unbußfertige Verfolger seiner Kirche, und tröste und errette alle gläubige Kreuzträger, zu Ehren seinem heiligen Namen, Amen.

\chapter{Das 13. Kapitel}

\section{Vorrede}

Wer Gott treulich von ganzem Herzen und rechtschaffen dienet, dem dienet Gott wieder auch nutzbarlich. Gott dienen aber heißt, Gottes Wort folgen und gehorsam sein, beides im Glauben gegen ihn und in der Liebe gegen den Nächsten. Derwegen weil wir Gottes auch keinen Augenblick ohne unsern Schaden entbehren können, und weil wir Gottes zu unsern Augen alle Augenblicke bedürfen: so laßt uns durch Hülfe des heiligen Geistes Gottes Willen thun, und ihm dienen ohne Furcht unser Leben lang in Heiligkeit und Gerechtigkeit, die ihm gefällig ist, Luc. 1. Und darauf höret an das 13. Kapitel.

\subsection{Abram scheidet sich von Lot, seinem Vetter}
\quellseite{59}

\section{Summarium}

Hält uns erstlich ein schön Exempel für mit Abraham und Lot, wie Gott seinen Segen reichlich mittheilet denen, so ihn fürchten, nach dem Spruch Christi: Trachtet am ersten nach dem Reich Gottes und nach seiner Gerechtigkeit, so wird euch das andere alles zufallen; denn euer Vater weiß, was ihr nothdürftig seid, Matth. 6, 33.

Darnach wie Uneinigkeit zu fliehen und sonderlich die Blutfreunde, alles, so Widerwillen erregen kann, meiden und eher den Vortheil aus den Händen geben sollen, denn daß man Ursache zu Zank und Widerwillen geben sollte, wie der fromme Abraham mit Lot thut. Solcher Fleiß ist sonderlich vonnöthen; denn wer Zank suchet und nach Einigkeit nicht mit besonderm Fleiß trachtet, der wird sein allenthalben genug finden. Aber selig sind die Friedfertigen.

\section{Votum}

Gott der Allmächtige steure in aller Welt dem Zanke und Hader und allen andern Sünden, und regiere uns, daß wir leben nach seinem Willen, ihm zu Ehren, uns zu Nutz und zu unsers Nächsten Wohlfahrt, Amen.

\chapter{Das 14. Kapitel}

\section{Vorrede}

Gott züchtiget mit väterlicher Ruthe Land und Leute, daß sie sollen Buße thun, auf daß er sie nicht mit dem Garaus und ewiger Pein strafen dürfe. Das bezeuget St. Paulus 1 Cor. am 11. und saget: Wenn wir gerichtet, das ist, gestrafet werden, so werden wir von dem Herrn gezüchtiget, auf daß wir nicht sammt der Welt verdammt werden. Derowegen weil Gott mit allerlei Krankheiten, mit seltsamen Gewittern, mit wunderlichen Fällen und auf andere mancherlei Weise uns zur Buße rufet, so wollen wir unsere Sünde erkennen, die Gott abbitten in wahrem Glauben an den Herrn Christum, und wollen gottsfürchtig leben, daß uns Gott mit endlichem Verderben nicht strafen dürfe.

Darauf höret an das 14. Capitel.

\subsection{Abram errettet Lot; wird von Melchisedech gesegnet}
\quellseite{60}

\section{Summarium}

\unsicher{Sie} gehet an die Legende des heiligen Patriarchen Abrahams, welchen Gott fordert aus dem Lande Haran ins Land Canaan, und verheißet ihm, er wolle ihn zum besondern Volk machen, und das Land Canaan eingeben; und zu seiner Zeit ihm einen solchen Samen erwecken, der ein gebenedeyter Same sei; das ist, durch welchen alle Welt aus dem Fluch zum Segen und Gottes Gnade soll gebracht werden. Sie soll man erstlich lernen, daß Gott gnädig sei. Darnach soll man Acht haben auf das Exempel Abrahams, wie er Gott trauet und seinem Befehle folget und Gottes Wort prediget, und dennoch das Kreuz nicht außen bleibt, finemal er Theurung halben in Egypten fliehen und da seines Weibes halben in Gefahr kommen muß. Aber weil er Gott glaubet und vertrauet, so hilft er ihm aus aller Noth; wie durch seine ganze Legende zu sehen ist; und strafet, die ihn beleidigen, wie er hie Pharao und sein Haus plaget, darum, daß er Sara ihm zum Weibe wollte nehmen.

\section[Notum]{\unsicher{Notum}}

Der gerechte starke Gott strafe noch heut alle muthwillige und unbußfertige Verfolger seiner Kirche, und tröste und errette alle gläubige Kreuzträger, zu Ehren seinem heiligen Namen, Amen.

\chapter{Das 13. Kapitel}

\section{Vorrede}

Wer Gott treulich von ganzem Herzen und rechtschaffen dienet, dem dienet Gott wieder auch nutzbarlich. Gott dienen aber heißt, Gottes Wort folgen und gehorsam sein, beides im Glauben gegen ihn und in der Liebe gegen den Nächsten. Derowegen weil wir Gottes auch keinen Augenblick ohne unsern Schaden entbehren können, und weil wir Gottes zu unserm Nutzen alle Augenblicke bedürfen: so laßt uns durch Hülfe des heiligen Geistes Gottes Willen thun, und ihm dienen ohne Furcht unser Leben lang in Heiligkeit und Gerechtigkeit, die ihm gefällig ist, Luc. 1. Und darauf hört an das 13. Kapitel.

\subsection{Abram scheidet sich von Lot, seinem Vetter}
\quellseite{61}

\section{Summarium}

Hält uns erstlich ein schön Exempel für mit Abraham und Lot, wie Gott seinen Segen reichlich mittheilet denen, so ihn fürchten, nach dem Spruch Christi: Trachtet am ersten nach dem Reich Gottes und nach seiner Gerechtigkeit, so wird euch das andere alles zufallen; denn euer Vater weiß, was ihr nothdürftig seid, Matth. 6, 33.

Darnach wie Uneinigkeit zu fliehen und sonderlich die Blutfreunde, alles, so Widerwillen erregen kann, meiden und eher den Vortheil aus den Händen geben sollen, denn daß man Ursache zu Zank und Widerwillen geben sollte, wie der fromme Abraham mit Lot thut. Solcher Fleiß ist sonderlich vonnöthen; denn wer Zank suchet und nach Einigkeit nicht mit besonderm Fleiß trachtet, der wird sein allenthalben genug finden. Aber selig sind die Friedfertigen.

\section{Votum}

Gott der Allmächtige steure in aller Welt dem Zanke und Hader und allen andern Sünden, und regiere uns, daß wir leben nach seinem Willen, ihm zu Ehren, uns zu Nutz und zu unsers Nächsten Wohlfahrt, Amen.

\chapter{Das 14. Kapitel}

\section{Vorrede}

Gott züchtiget mit väterlicher Ruthe Land und Leute, daß sie sollen Buße thun, auf daß er sie nicht mit dem Garaus und ewiger Pein strafen dürfe. Das bezeuget St. Paulus 1 Cor. am 11. und saget: Wenn wir gerichtet, das ist, gestrafet werden, so werden wir von dem Herrn gezüchtiget, auf daß wir nicht sammt der Welt verdammt werden. Derowegen weil Gott mit allerlei Krankheiten, mit seltsamen Gewittern, mit wunderlichen Fällen und auf andere mancherlei Weise uns zur Buße rufet, so wollen wir unsere Sünde erkennen, die Gott abbitten in wahrem Glauben an den Herrn Christum, und wollen gottfürchtig leben, daß uns Gott mit endlichem Verderben nicht strafen dürfe. Darauf hört an das 14. Kapitel.

\subsection{Abram rettet Lot; wird von Melchisedek gesegnet}
\quellseite{62}

\textit{Zu Vers 18.}a Trug Brod) nicht daß ers opferte, sondern daß er die Gäste speisete und ehrte, dadurch Christus bedeutet ist, der die Welt mit dem Evangelio speiset.

\section{Summarium}

Dies ist eine Historia, wie Gott durch leibliche Strafe endlich die Menschen zur Buße fordert, daß sie sich bessern und frömmer sollen werden. Denn ehe denn er Sodoma und Gomorra ließ mit höllischem Feuer verbrennen, vermahnet er sie zur Buße und Besserung durch Kriege.

Solche Strafe muß der fromme Lot auch mit tragen. Aber endlich wird er errettet durch Abraham, welcher hie zum Exempel stehet aller Obrigkeit, daß sie die Ihren schützen und doch nicht eigennützig sein sollen.

Daß Abraham dem König Melchisedek den Zehnten gibt, welcher ein Priester Gottes war, ist zum ersten ein Exempel, daß weltliche Obrigkeit schuldig ist zu helfen, daß die Kirchen und der rechte Gottesdienst erhalten werde. Denn sie muß Schutz und Unterhaltung haben. Und selig sind die, so ihr Geld, Gut und Gewalt zu solchem köstlichen guten Werke brauchen, daß das Predigtamt dadurch gefördert und erhalten werde.

Danach hat es seine Deutung auf das Priesterthum Christi, wider das Gesetz, wie die Epistel an die Ebräer zeuget, Kap. 7, 10. und David Ps. 110, 4., der Christum einen Priester nennet, nicht nach der Ordnung oder Weise Levi, sondern Melchisedek, der nicht mit Opfern und Sünden zu thun hat, wie der Priester im Gesetz, sondern segnet und nimmt den Fluch von uns. Der Name Melchisedek heißt auf deutsch so viel, als ein König der Gerechtigkeit. Solchen Namen führet unser lieber HErr Christus, der durch Vergebung der Sünden uns gerecht machet. Es ist aber gewiß dieser Melchisedek Sem gewesen, des Noah Sohn, welcher dazumal der älteste Patriarch und der oberste Regent und Lehrer in der Kirche gewesen ist, und derhalben den Namen hat, er sei des Allerhöchsten Priester. Dieses Sems Nachkommen ist Abraham, und hat Abraham von ihm geerbet die Verheißung von Christo.

\section{Votum}

Der treue barmherzige GOtt verleihe, daß wir alle in wahrer Buße leben und sterben; und segne und behüte ferner unsere Stadt, und das ganze gemeine Vaterland, um Jesu Christi unsers HErrn willen, Amen.

\chapter{Das 15. Kapitel}

\section{Vorrede}

Die großen Heiligen haben auch großes Kreuz und müssen eines nach dem andern tragen. Aber sie sind davon, auch mit großer Freude, endlich erlediget worden. Von beiden sagt der Prophet David im 34. Psalm also: Der Gerechte muß viel leiden, aber der HErr hilft ihm aus dem allem. Derowegen wer ein Kreuzträger ist, der lasse ihm Gottes guten väterlichen Willen wohlgefallen, sei geduldig und bete. Denn Gott

2*
\quellseite{63}
zürnet nicht ewiglich und seine rechte Hand kann alles ändern, Ps. 77.; und kann auch das größte Unglück bald in Freude verwandeln. Dieses lehret auch das 15. Kapitel.

\subsection{Abrams Glaube und Gerechtigkeit wird gepriesen}

\textit{Zu Vers 9.}a Bringopfer.

\textit{Zu Vers 11.}b Gevögel) das Gevögel und der rauchende Ofen und der feurige Brand bedeuten die Egypter, die Abrahams Kinder verfolgen sollten; aber Abraham scheucht sie davon, das ist, Gott erlöset sie um der Verheißung willen, Abraham versprochen. Daß aber er nach der Sonne Untergang erschrickt, bedeutet, daß Gott seinen Samen eine Zeit lang verlassen wolle, daß sie verfolget würden, wie der HErr selbst hie betet. Also gehet es auch allem Gläubigen, daß sie verlassen und doch erlöset werden.

\section{Summarium}

Zeiget erstlich an, wie der fromme Abraham ohne große Anfechtung nicht sei gewesen, daß Gott die Verheißung, daß er ihn zum großen Volk wollte machen, so lange verzog. Darum, da Gott zu ihm spricht: Fürchte dich nicht, ich bin dein Schild etc.; spricht er darauf: Was willst du mir geben, ich gehe dahin ohne Kinder; als wollte er sagen: Soll ich zum großen Volk werden, so mußt du mir, lieber HErr, ja einen Sohn geben. Als ihm Gott aber verheißet, er soll einen Sohn haben, überwindet er solche Anfechtung, daß er so alt und seine Sara unfruchtbar war, und glaubet Gott. Solcher Glaube, spricht Moses, ist ihm zur Gerechtigkeit gerechnet. Darnach am Ende ist eine Weissagung, wie die Verheißung aufgezogen und Abrahams Nachkommen zuvor in Egypten müssen geplagt werden, auf daß sie sich an solchem Unfall nicht ärgerten, sondern fest an der Zusagung Gottes hielten, daß Land Canaan würde noch ihr eigen werden, und daß Christus, der gebenedeite Same, Gottes Verheißung nach, gewißlich kommen würde.

\section{Psalm}

Wir danken GOtt, der seiner Christenheit unter dem Kreuz beigestanden ist allezeit;
\quellseite{64}
der wolle uns auch bewohnen, und in aller Noth trösten und erhalten, hochgelobt in Ewigkeit, Amen.

\chapter{Das 16. Kapitel}

\section{Vorrede}

Wenn ein böses Ding ein allgemeiner Gebrauch wird, so pflegen solche auch die, so fromme Leute sonst sind, nicht für so gar sträflich zu halten: wie heute viele große Laster, als Trunkenheit, Hoffart, Wucher und dergleichen gering geachtet werden, weil die beim größten Haufen im Schwange gehen. Aber Sünde ist und bleibt Sünde, und wenn sie noch so gemein wäre. Derowegen wer also handeln und wandeln will, daß er gegen Gott bestehe, der soll nicht nach der Welt bösen Exempeln, sondern nach Gottes ernstem Willen und unwandelbarem Wort leben. Denn das allein soll unser Fuß Leuchte sein und ein Licht auf unserm Wege, saget der königliche Prophet David Ps. 119. Dazu höret auch an das 16. Kapitel.

\subsection{Hagar gebieret dem Ismael}

\textit{Zu Vers 2.}a Bauen) das ist, Kinder kriegen.

\textit{Zu Vers 11.}b Ismael) heißt, Gott erhöret.

\section{Summarium}

Zum ersten ist ein Exempel, wie die Vernunft Rath suchet, ihr selbst zu helfen, wo Gott die Hülffe zu lange verzeucht und nicht so bald hilft, als man hofft. Denn Sarai, weil sie je länger je älter wird, und doch unfruchtbar bleibt, beginnt sie an Gottes Verheißung zu wanken, denket, es werde mit ihr nichts draus, weil sie unfruchtbar und nun alt sei, und hält an bei Abram, daß er die Hagar, ihre Magd, zum Weibe nehme, auf daß Gottes Verheißung nicht verhindert werde. Es hat aber Gott solches lassen also geschehen zur sondern Deutung; wie Paulus zu den Galatern Kap. 4, 22. davon disputirt, von zweien Kirchen, die eine so mit Werken umgehet und auf eigen Verdienst bauet; die andere, so an Gottes Zusagung und Barmherzigkeit sich hält.

Darnach ist hie eine feine Hauslehre, wie Hagar, da sie schwanger war, ihrer Frau Sarai Strafe nicht leiden wollte und lief davon. Aber der Engel heißt sie wieder in den Dienst gehen und ihrer Frau unterthan sein, und sich für ihr demüthigen; zum Exempel allen Ehehalten oder Dienstboten, daß sie mit Demuth allen Gehorsam ihrer Herrschaft leisten sollen. Denn sol-
\quellseite{65}
Dieß will Gott haben und gefällt ihm wohl, will auch Glück dazu geben; so dagegen, wo das Gesinde untreu und ungehorsam ist, Gott es mit allerlei Unglück strafen will.

\section{Datum}

Verleihe der gütige treue Gott, daß wir uns nach seinem Wort gegen ihn und gegen unsern Nächsten gehorsamlich verhalten, und daß wir an Leib und Seele seinen Segen haben und behalten, Amen.

\chapter{Das 17. Kapitel}

\section{Vorrede}

Der Glaube und der Gehorsam gegen Gott sind die zwei Stücke, welche die heilige Schrift fürnehmlich lehret, und auch von allen Menschen, die zeitliche Wohlfahrt und ewige Seligkeit erlangen sollen, erfordert. Derowegen läßt uns alle Tage fleißig zur Kirche gehen und beten, daß die Anhörung göttlichen Wortes uns auch zu täglicher Stärkung unsers Glaubens und Besserung unsers Lebens fruchtbarlich diene und gelange. Also wird sich Gott gegen uns väterlich erzeigen, da wir uns gehorsamlich gegen sein Wort verhalten. Darauf höret auch an das 17. Kapitel.

\subsection{Isaaks Verheißung wird mit der Beschneidung bestätigt}

\textit{Zu Vers 5.}a Abram heißt hoher Vater, Abraham aber, der Haufen Vater; wiewohl dieselben Haufen nur mit einem Buchstaben angezeigt werden in seinem Namen, nicht ohne Ursache.

\textit{Zu Vers 7.}b Nachkommen) Ledorotham, daß ist, so lange ihr Ding währen wird. Denn Moses hiermit deutet, daß ihr Ding solle endlich aufhören und ein anderes kommen.

\section{Summarium}

Ist eine Historia, wie Gott Abraham tröstet und die Verheißung mit einem neuen Bundeszeichen der Beschneidung bekräftiget, daß er aller der Gott sein, sie segnen und behüten wolle, so solches Bundeszeichen an ihrem Leibe tragen.

Die hebt sich erstlich der Unterschied an, daß Gott den Abraham und seine Nachkommen als ein sonder Volk von allen andern Völkern absondert und ihn und seine Nachkommen, die Juden, für der Welt und an ihrem Leibe mit einem sondern Gemerk oder Zeichen bezeichnet; auf daß man wüßte, daß Christus allein aus diesem beschnittenen Volk und sonst von keinem Volk sollte geboren werden.

Zum andern hat solche Beschneidung auch ihre Bedeutung gehabt: denn Gott hat damit wollen anzeigen, wie unsere Natur sündhaftig sei und dürfe, daß sie geistlich beschnitten würde; wie die Propheten auch predigen, daß man die Vorhaut des Herzens beschneiden, das ist, in Gottesfurcht wandeln und die Sünde mit allen bösen Lüsten töten soll.

Zum dritten ist diese Beschneidung gewesen,
\quellseite{66}
wie es Paulus Röm. 4, 11. nennet, ein Siegel der Gerechtigkeit; denn dabei hat Abraham und alle seine Nachkommen sich können erinnern, daß sie Gottes Volk sind; sintemal solches Bundeszeichen, als ein Sacrament, die Verheißung hätte. Gott wollte solcher beschnittenen Leute sich sonderlich annehmen, sie schützen und segnen als die Seinen.

\section{Summarium}

Ist endlich ein Exempel, daß wir lernen sollen, daß Gottes Verheißungen wahr sind und wir daran nicht sollen zweifeln; denn bei Gott ist nichts unmöglich; wie es Abraham unmöglich däucht, daß er von seiner Sarah noch sollte einen Sohn zeugen.

Darnach ist's ein Exempel eines christlichen Gehorsams in Abraham, daß er nicht lange das von disputiret, warum er und sein Gesinde sich beschneiden sollen; sondern weil es Gott befohlen hätte, thut ers und folget mit einfältigem Glauben, B. 23. Denn es für der Welt ein sehr närrisch Ansehen hat, daß der alte Mann sich und sein ganzes Haus beschneiden soll. So kann Vernunft keinen sondern Nutzen noch Frucht solches Werkes erdenken. Aber um solches bekümmert sich Abraham gar nichts und gehet sein straß in vollem Gehorsam und willigem Herzen dahin, läßt ihm an dem genügen, daß es Gott also geheißen und befohlen hat. Also sollen wir in allen Dingen auch thun und nur dahin sehen, ob es Gottes Befehl sei, und weiter nichts fragen noch uns bekümmern, warum ers befohlen oder wozu es nütze sei. Solchen einfältigen Gehorsam will Gott von uns allen haben.

\section{Votum}

Der barmherzige gütige Gott verleihe, daß wir seinem Willen in Liebe und Leid treulich folgen, und seiner gnadenreichen Verheißung zu zeitlicher und ewiger Wohlfahrt fruchtbarlich genießen, Amen.

\chapter{Das 18. Kapitel}

\section{Vorrede}

Gottes Wort hat diese Kraft, wo man es mit Ernst handelt, höret und betrachtet, daß es nimmer ohne Frucht abgehet, sondern allzeit neuen Verstand und Andacht erweckt und reine Herzen und Gedanken machet. Auch ist dies gewiß: wie sich der Mensch gegen das Wort Gottes verhält in seinem Leben, also verhält sich Gottes Wort wiederum gegen den Menschen in seinem Absterben. Derowegen wollen wir Gottes Wort lieben und lernen, und bitten, daß Gott auch dadurch in uns täglich wirken wolle Verstand, Glauben und Gehorsam, und daß wir daraus Freude und Trost haben unterm Kreuz, und in unserer Hinfahrt aus diesem zeitlichen kümmerlichen Leben, Amen. Darauf höret ferner an das 18. Kapitel.

\subsection[Dem Abraham wird nochmals Isaac verheißen, und die]{Dem Abraham wird nochmals Isaac verheißen, und die Verheißung Sodoms geoffenbaret}
\quellseite{67}

\textit{Zu Vers 2.}a Nieder] vor einem fällt er nieder und redet auch als mit einem, und doch mit dreien. Da ist die Dreifaltigkeit von Gott angezeigt.

\textit{Zu Vers 14.}b So ich lebe] Gott redet als ein Mensch, 1 Mos. 3, 9.: Adam, wo bist du? Item 1 Mos. 11, 7.: Ich will hinab fahren und sehen. 1 Mos. 18, 21.: Ich will sehen, ob sie sei. Denn dies Wort (so ich lebe), setze es, wo du willst, so lautet es doch nicht, als wenn er von der Frucht oder Eltern wollte sagen: Du sollst einen Sohn haben, so er lebet, oder, so ihr lebet. Meinst du, er wisse nicht, ob er Sohn oder die Eltern leben werden? Ja, wo sie nicht ewig leben würden, was wäre die Verheißung?

\section{Summarium}

Ist eine holdselige Historia, wie Gott zu Abraham kommet und sein Gast wird, um \unsicher{meinen} darum, daß er ihn im Glauben stärke, daß er nicht zweifele, er werde einen Sohn von seiner Sarah haben, denn dem HErrn sei nichts unmöglich, V. 10. Auf daß wir auch lernen auf Gottes Wort sehen, und uns gewiß auf seine Zusagung verlassen. Sonderlich aber ist hie zu merken, wie der Text sagt, Abraham habe drei Männer gesehen, und redet doch je als mit einem allein, und je als mit dreien, V. 2.

Item, der HErr sprach zu Abraham etc., V. 13.

Solches ist eine gewisse Anzeigung, daß nur ein einiger Gott sei, aber drei unterschiedliche Personen, als GOtt Vater, GOtt Sohn, GOtt heiliger Geist. Wie denn die Offenbarung im Neuen Testament solches klärlich mit bringt. Solche Zeugnisse aber des Alten Testaments dienen dazu, daß wir lernen, die Patriarchen haben auch solches gegläubet und gewußt wie wir, und sind eben in solchem Glauben selig worden.

Darnach ist ein fein Exempel Abrahams halben, der so gastfrei und gegen fremde und unbekannte Gäste so freundlich und diensthaftig ist, V. 4, nach der Regel St. Pauli Röm. 12, 13.: Herberget gerne.
\quellseite{68}

\textit{Zu Vers 25.}c In Hebräisch: sollst du sonst alle Welt richten und nicht recht thun.

\section{Summarium}

Ist erstlich eine köstliche Historia, in welcher Gott seine Barmherzigkeit reichlich läßt sehen, wie er gerne alle Strafe abwenden und schonen wolle, so nur die Menschen sich begehreten zu bessern; ja daß er um wenig frommer Leute willen eines ganzen Landes gerne verschonen wollte.

Widerum aber ist es schrecklich, daß die Welt so arg und boshaftig soll sein, daß unter so viel Leuten in einem ganzen Lande nicht zehn fromme und gottesfürchtige Menschen sollen erfunden werden. Darum muß die Strafe über sie gehen und Abrahams Fürbitte umsonst sein, weil keine Besserung des Lebens folget. Denn beten und sich bessern muß beides bei einander sein.

Zum andern siehet man hie, wie die Fürbitte für andere gegen Gott so ein trefflich Ding ist. Aber wer solcher Fürbitte genießen will, muß sich dargegen recht schicken und die Fürbitte durch Unbußfertigkeit nicht hindern.

\section{Gebet}

Der allmächtige barmherzige Gott komme auch zu uns, und verleihe, daß wir glauben, beten und leben, daß sein Name allzeit durch uns geheiliget, gelobet und gepreiset werde, hie zeitlich und dort ewiglich, Amen.

\chapter{Das 19. Kapitel}

\section{Vorrede}

Betet herzlich, daß Gott der heilige Geist in uns kräftig wirken wolle, auf daß die erschreckliche Historia vom Untergang Sodoma und Gomorra heut so tief in unsere Herzen sich senke, daß wir der nimmermehr vergessen, und auf daß wir diese Geschichte zu unserer täglichen Besserung fruchtbarlich gebrauchen. Denn Gott dem HErrn niemand so lieb ist, da er sündigt, und in Sünden muthwillig verharrret, er straft ihn, daß ers an Leib und Seele, zeitlich und ewig fühlet. Aber, o besser nicht geboren, denn allhie Gottes Zorn tragen, und dort ewig sein verloren. Darauf höret an das 19. Kapitel.

\subsection{Vertilgung Sodoma. Lots Errettung und Sündenfall}
\quellseite{69}

\textit{Zu Vers 22.}a Zoar) heißt klein.
\quellseite{70}

\section{Summarium}

In dieser Historie stehet man, wie Gott an seine Heiligen denkt, und sie errettet von der Strafe, so über die Gottlosen kömmt. Denn die Frommen sollen ihrer Frömmigkeit genießen, wie Lot.

Wiederum stehet man, weil die Sünde so überhand genommen hat, und die Sodomiter keine Warnung wollten zu Herzen nehmen, daß Gott hat müssen solchen Ernst brauchen, auf daß wir lernen uns fürchten und für Sünden hüten; sintemal sie Gott nicht kann ungestraft hingehen lassen. Merke aber wie sicher die Sünder sind, wenn ihr Verderben am allernächsten ist: die Sünde des Lots, da sie gewarnet werden, hielten es für ein Gelächter und gedachten nicht zu weichen, B. 14.

Daß die Frau Lots sich umsiehet und zur Salzsäule wird, B. 26. deutet Christus im Evangelio Luc. 17, 21. dahin, daß man sich das Zeitliche nicht lasse hindern, sondern auf Gottes Wort sehe und demselbigen folge; sintemal wir der Strafe sonst nicht können entrinnen. Solche Deutung nimmt der Herr gewöhnlich davon, daß es mit Lots Weib nicht ein schlechtes fürwitziges hinter sich sehen gewesen ist, sondern daß sie ihr die Welt auch hat lassen lieb sein, und Gottes ernstliches Dräuen nicht für wahr gehalten hat. Solchen Unglauben hat Gott mit Ernst strafen, und das Gedächtnis andern zur Abscheu, öffentlich wollen bleiben lassen.

\section{Summarium}

Ist eine Historia, welche nicht darum geschrieben ist, noch gelesen wird, daß wir solchem Exempel, Lots und seiner Töchter, folgen, sondern unsere eigene Bosheit recht sollen erkennen lernen, und desto mehr Gott für Augen haben, und unsere bösen Lüste dämpfen; sintemal auch die Heiligen sich nicht allezeit bewahren können, sondern bisweilen allzu grob fallen.

Der Text zeiget an, daß die Töchter Lot solches aus Einfalt fürgenommen, und Lot unwissend gethan habe, B. 31. 33. Aber damit sind sie beide nicht entschuldiget. Und sonderlich stehet man hie, was für gräuliche Sünden aus der Trunkenheit erwachsen: denn wo Lot nüchtern gewesen, würde er nimmermehr solche Sünde begangen haben. Ob er nun wohl als ein trauriger, bekümmerter Mann, der so vieler Seelen Verderben gesehen, so viel guter Freunde und unter denselben sein liebes Weib alle auf einen Tag, und schier in einem Augenblicke verloren hatte, mehr vor Leid und Kummer, denn von Wein ist voll worden; denn bekümmerten Leuten schadet der Trunk viel leichter: so ist die Trunkenheit an ihr selbst ein schädliches Laster, und verursacht viel Jammer und Noth; darum alle Christen sich dafür hüten sollen.

\section{Gebet}

Helffe der treue gütige GOtt, daß alle Welt an Sodoma und Gomorra sich spiegele, Buße thue, und Gottes Zorn und Strafe in diesem und künftigen Leben entfliehe, Amen.

\chapter{Das 20. Capitel}

\section{Vorrede}

Keine Sünde ist so groß, gräulich und schädlich, ein jeder unter uns würde solche schon lange begangen und gethan haben, da uns Gott nicht aus väterlicher Güte und Barmherzigkeit
\quellseite{71}
behütet hätte. Dafür laßt uns ihm herzlich danken und bitten, er wolle uns ferner regieren und führen, daß wir durch unsere Sünde ihn nicht erzürnen, des Herrn Christi nicht spotten, den heiligen Geist nicht betrüben, die lieben Engel nicht verjagen, uns und den Unsern, auch unserm lieben Vaterlande, verderbliche Strafen nicht zuziehen; und laßt uns derhalben täglich mit dem Könige David seufzen, Psalm 119. Thue wohl deinem Knecht, daß ich lebe, und dein Wort halte; also: thue wohl deiner Magd, daß ich lebe und dein Wort halte. Zu diesem höret an das 20. Kapitel.

\subsection[Sarah wird dem Abraham von Abimelech genommen, und]{Sarah wird dem Abraham von Abimelech genommen, und mit Gewinst wieder gegeben}

\textit{Zu Vers 10.}a Geschehen weil du ein Prophet bist, magst du was gesehen haben, daß ichs verdienet habe mit meinen Sünden.

\textit{Zu Vers 13.}b Wandern Gott hieße mich in die Irre ziehen, als wären viel, und doch ein Gott.

\textit{Zu Vers 16.}c Ihre Strafe hieß die Heiligen werden säuberlich und mit Gewinst gestraft, als hie Sarah wird gestraft, daß sie Abraham hatte Bruder genennet, und erhält große Wohlthat.

\section{Summarium}

Ist eine Historia, wie Abraham seiner Sarah halben noch einmal in Gefahr kommt, daß sie ihm vom König Abimelech genommen wird, V. 2. Aber hie beweiset sich Gott wiederum, wie er die Seinen behüten wolle, und sonderlich wie er dem Ehebruch feind, und denselben mit Unfruchtbarkeit und dem Tod wolle strafen, V. 18. Daneben ist sonderlich wohl zu merken, wie es eine Gabe Gottes und von Gott zu bitten sei, daß er für Sünden uns bewahren wolle; wie Gott zu Abimelech sagt: Ich habe dich behütet, daß du nicht wider mich sündigest,
\quellseite{72}
und habe das nicht zugeben, daß du Sarai, Abrahams Weib, berühretest, V. 6. Denn unsere Natur ist viel zu schwach, daß sie sich aller Anfechtung und Sünde erwehren sollte, wo Gott nicht mit sonderer Gnade hält und hütet, und wir nicht täglich um solchen Schutz und Hülfe bitten; wie Christus im Vaterunser auch befiehlt und lehret, daß wir beten sollen: Nicht einführe uns in Versuchung, das ist, lieber Herr Gott, wo wir in Versuchung und Anfechtung geraten, sei gnädig, und führe uns wieder heraus, daß wir nicht darinnen liegen bleiben und verderben. Solcher Bitte bedürfen wir überaus wohl.

\section{Votum}

Gott der heilige Geist helfe, daß wir in täglichem Gebet fleißig und andächtig sein; und stärke uns an Leib und Seele, daß wir recht gläubig und christlich leben, bis an unser selig Ende, Amen.

\chapter{Das 21. Kapitel}

\section{Vorrede}

Freilich ist die heilige Biblia wie ein großer schöner Garten, darin viel und allerlei Arten Bäume stehen, davon man auch allerlei gut Obst und Früchte haben kann, die zur Speise, Arznei und Labsal dienen. Denn die heilige Biblia gibt fleißigen Zuhörern und andächtigen Herzen reichen Unterricht, Vermahnung, Warnung, Verheißung und Trost, die nützlich sein in Leib und Leid, Noth und Tod. Derowegen laßt uns die heilige Biblia also anhören und lernen, auf daß wir täglich zu unserer Besserung, und zu unserm Trost, etwas mit zu Hause nehmen, allezeit behalten, fruchtbarlich im Leben und Sterben gebrauchen. Folgendes höret an das 21. Kapitel.

\subsection{Isaaks Geburt. Austreibung Ismaels. Abrahams Bund mit Abimelech}

\textit{Zu Vers 17.}a Hagar) mercke hier auf Hagar, wie die des Gesetzes und glaubensloser Werck Figur ist, Gal. 4, 22, und dennoch sie Gott zeitlich belohnet und groß machet auf Erden.
\quellseite{73}

\section{Summarium}

Ist eine Historia, wie Sarah in ihrem Alter einen Sohn träget. Auf daß wir sehen, wie bei Gott nichts unmöglich und sein Wort allein wahr sei, und wir uns desto gewisser darauf verlassen.

Hie hebet sich nun die Spaltung der Kirchen: Ismael bedeutet, wie St. Paulus Gal. 4, 21. sagt, das Volck, so durch die Werke des Gesetzes will selig werden. Aber da wird nichts aus: der Ismael muß aus dem Haus mit seiner Mutter und sich am zeitlichen Segen begnügen lassen, Isaac aber bleibet im Hause und ererbet den Segen, B. 10. Das sind die, so nicht auf ihre Werke, sondern auf Gottes Güte und Barmherzigkeit ihr Vertrauen setzen.

Ueber das ist hie auch ein fein Exempel, wie Gott Ismaels Gebet erhöret, der nun ein Knabe bei vierzehn Jahren alt war und sein Leben rettet, B. 17.; auf daß die jungen Leute auch lernen, mit Ernst beten und fromm sein. Denn Gott will gerne sie schützen und in aller Noth helfen.

\textit{Zu Vers 31.}b Versaba) heißt auf deutsch Schwörbrunn oder Eidsbrunn, möchte auch wohl Siebenbrunn heißen.

\section{Summarium}

Am Ende dieses Kapitels ist ein Exempel, daß man in bürgerlichen Sachen der Einigkeit zu gut, ohne Sünde möge Eid schwören, wie Abraham und Abimelech thun, B. 23. Darnach stehet, Abraham habe geprediget vom Namen des HErrn, des einigen Gottes, B. 34., das ist, von der Verheißung von Christo: wie Gott durch Christum in unserm Fleisch sich offenbaren und ein ewiges Reich werde anrichten. Darum sagt Christus im Evangelio, Abraham habe seinen Tag gesehen und sich gefreuet.

\section[Notum]{\unsicher{Notum}}

GOtt der HErr erleuchte uns, daß wir uns auch in wahrem Glauben unsers HErrn JEsu Christi freuen; in unsern Nöthen fleißig beten, und auf sein Wort uns beständig verlassen, bis zum ewigen Leben, Amen.

\chapter{Das 22. Kapitel}

\section{Vorrede}

Eines jeden Christen Menschen Leben soll eine tägliche Buße sein. Dieses lehret der heilige Apostel Paulus, da er sagt, 1 Tim. 1.: Uebe eine gute Ritterschaft, habe den Glauben und gut Gewissen. Derowegen wollen wir heute abermal unsere Sünde erkennen und sonderlich diese, daß wir alle schwach im Glauben sind, und daß wir noch nie Gott dem HErrn den Gehorsam geleistet haben, den wir ihm schuldig sind. Und wollen ihn herzlich bitten, daß er uns, aus Gnaden um Christi willen, alle unsere Sünde vergebe, unsern Glauben uns vermehre und stärke und uns erfülle mit Früchten der Gerech-
\quellseite{74}
tigkeit, auf daß wir sind lauter und unanstößig bis auf den Tag Christi unsers Herrn, Amen. Daneben höret an das 22. Kapitel.

\subsection{Aufopferung Isaaks. Verheißung von Christo}

\textit{Zu Vers 2.}a Moria) Moria heißt Gottesfurcht, Reverentia Dei, Cultus Dei. Denn die Altväter Adam, Noah, Sem, auf demselbigen Berge, Gott gehret, gefürchtet, gedienet haben. Wir Deutschen heißen es vielleicht den heiligen Berg, oder da man Gott dienet mit Loben, Beten und Danken.

\textit{Zu Vers 8.}b Ersehen) Gott siehet und weiß wohl, wo das Schaf sei, läßt ihn dafür sorgen, er siehet es besser, denn wir.

\textit{Zu Vers 14.}c Siehet) Hebrei dicunt, Dominus videbitur. Sed nos Hieronymum secuti Rabinos Grammaticos cum suis punctis et Cametz hoc loco negligimus, et sine punctis dicimus: Der HErr siehet, das ist, Gott sorget für alles und machet. Etiamsi sensus ille: Dominus videbitur, sit pius valde, quod Deus appareat, ubi verbum ejus docetur, quod Rabini Grammatici non intelligunt, Heb. 6, 14.

\textit{Zu Vers 18.}d Gehorschet) hie wird Abraham nicht gerecht durch seinen Glauben für sich, sondern erweiset solche Herrlichkeit seines Samens; denn er zuvor gerecht ist, wie oben Cap. 15, 6.
\quellseite{75}

\section{Summarium}

Dies ist ein schön Exempel eines sonderlichen großen Glaubens und völligen Gehorsams gegen Gott, daß Abraham Gottes Gebot folget und seines einigen Sohnes nicht verschonet; er will, wie Gott ihm befohlen hatte, ihn aufopfern, V. 6., so doch die Verheißung alles zeitlichen und geistlichen Segens allein auf dieses Sohnes Leben stund.

Auf daß wir auch lernen, uns ganz Gott ergeben und dennoch gewisse Hülfe erwarten; wie Abraham thut: läßt ihm seinen Sohn nicht zu lieb sein und denkt, Gott könne seine Zusage vollstrecken, wenn Isaac schon aufgeopfert werde. Solch gewiß Vertrauen auf Gottes Zusage ist Gottes höchster Dienst und Wohlgefallen.

Merke aber, daß der Text sagt: Gott versucht Abraham, V. 1., das ist, es war keine Meinung nicht: daß Abraham sollte seinen Sohn tödten. Denn Todtschlag ist verboten; er wollte nur dabei Abrahams Glauben, Gehorsam und Furcht erfahren.

\section{Votum}

Gott regiere uns alle, wie er den heiligen Abraham geführet hat; auf daß wir auch unser Licht leuchten lassen, daß die Menschen unsere gute Werke sehen, und Gott allezeit gelobet und gepreiset werde, Amen.

\chapter{Das 23. Kapitel}

\section{Vorrede}

Der Mensch lebet, und, wenn es ihm Gott noch so lange gönnet, so muß er doch endlich sterben. Wohl aber allen, die im Herrn entschlafen. Derowegen weil wir alle leben und wissen nicht, wie lange, und weil wir alle sterben und wissen nicht, wann; so wollen wir alle täglich Buße thun und also glauben, beten, handeln und wandeln, ob wir wohl zeitlich sterben müssen, daß wir doch mit Gott ewig leben, um Jesu Christi, unseres Heilandes, fröhlichen Geburt, bittern Leidens, schmählichen Todes, siegreichen Auferstehung und herzquickenden Himmelfahrt willen. Darauf höret nun ferner an das 23. Kapitel.

\subsection{Sarahs Tod und Begräbniß}

\textit{Zu Vers 2.}a Hebron) Hebron ist Kirjath Arba (spricht Moses), das ist, die Vierstadt; denn die hohen Hauptstädte waren vor Zeiten alle Arba, das ist, in vier Theile getheilet, wie Rom, Jerusalem und Babylon auch, 1 Mos. 10, 10.

\textit{Zu Vers 15.}b Sekel) Sekel ist ein Gewicht, an der Münze ein Drittheil gülden; denn vor Zeiten man das Geld so wog, wie man jetzt mit Gold thut.
\quellseite{76}

\section{Summarium}

Das ist eine Historia, wie Abraham seine liebe verstorbene Sarah beweint und beklaget und ihr ein eigen Begräbniß kaufet. Zum Zeugniß, daß er fest glaube, ob er wohl ein Fremdling im Lande sei, kein Theil darin habe, auch nicht eines Fußes breit, daß gleichwohl sein Same nach ihm dasselbe soll zu eigen haben.

Ist auch ein fein Exempel hie, wie im Kaufen und Verkaufen fromme Leute ehrbarlich und aufrichtig gehandelt haben. Wie solches noch wohl stehet, aber leider an solcher Ehrbarkeit täglich viel abgehet.

\section{Votum}

Helfe der treue GOtt, daß wir also leben und sterben, auf daß wir alle auferstehen zum ewigen Leben, das uns erworben hat JEsus Christus unser Heiland, hochgelobt mit dem Vater und heiligen Geist in Ewigkeit, Amen.

\chapter{Das 24. Kapitel}

\section{Vorrede}

Nach Gottes Wort ist der größte Schatz auf Erden, in einem gewissen und rechten Beruf stehen, und nichts anfangen oder fürnehmen ohne Gottes Befehl. Und wer derowegen seinen christlichen Beruf also treibet, daß er Gott um Segen bittet, ist fleißig und treu, und lebt in Gottesfurcht, dem gelinget es, daß er Gott dafür genug zu danken hat. Dieses lehret unser HErr JEsus Christus, da er sagt, Matth. am 7.: Bittet, so wird euch gegeben werden; und abermal, da er spricht: Trachtet am ersten nach dem Reich Gottes und nach seiner Gerechtigkeit, so wird euch das andere alles zufallen, Matth. am 6. Dessen ist ein Exempel das 24. Kapitel.

\subsection{Isaak heyrathet die Rebekka}
\quellseite{77}

\section{Summarium}

Ist eine Historia, in welcher wir sehen, wie die Eltern sich fleißigen sollen, daß ihre Kinder ehrlich und bei frommen Leuten heirathen. Und ist sehr tröstlich, daß Abraham und sein frommer Knecht Gott zum Heirathsmann oder Ehestifter nehmen und bitten, daß er den Isaac wohl versorgen wolle, B. 3. Und daß sie die Ehe so groß und ehrlich für Gott achten, daß Gott seine Engel dazu verordnen werde, daß es alles glücklich zugehe, B. 7. Dermaßen wird die Ehe bei den Papisten nicht geachtet, die es für einen solchen Stand achten, da Gott nichts mit zu thun habe, ja da man Gott nicht recht könne dienen.
\quellseite{78}

\textit{Zu Vers 47.}a Spangen) diese güldene Spange ist gewesen ein halber Zirkel auf der Stirn bis zu beiden Ohren; darum heißt ers jetzt Ohrenringe, jetzt Stirnpangen. Und siehet, als habens beide Mann und Weibsbilde getragen zum Schmuck. Wie unten Kap. 35, 4. Sprüchw. 11, 22. Circulus aureus in naribus suis. Das sagen wir Deutschen, die Sau gefrönet. Inde diadema Regum et lamina summi Sacerdotis in fronte etc.

\textit{Zu Vers 49.}b Saget mirs) er handelt zuvor mit Mutter und Bruder um die Braut. Daraus man siehet, daß heimliche Verlöbniß ohne Vorwissen der Eltern nicht recht ist.

\textit{Zu Vers 53.}c Würze) köstliche Früchte.

\textit{Zu Vers 57.}d Dirne) die Braut soll ungezwungen zur Ehe gegeben sein, von den Eltern, dazu auch gefragt werden um ihren Willen.
\quellseite{79}

\section{Summarium}

Ist ein fein Exempel, wie der Rebekka Freundschaft mehr auf Gottes Schickung stehet, denn auf Geld und Gut. Darum sind sie willig und lassen die Tochter in ein fremdes Land dem Isaak gerne folgen.

Wie sich nun der Bräutigam und der Werber dagegen halten, ist auch fein zu merken: Isaak gehet aufs Feld und betet, V. 63. Das soll ein Bräutigam thun, soll anders Glück und Heil, wie wohl es bedarf, dabei sein. So danket der Werber auch Gott, daß er seine Reise also gefördert und Glück dazu gegeben hat, V. 56. In Summa, du siehest, wie diese fromme, heilige Leute alles mit Gottesfurcht in rechtem Glauben und mit ernstem Gebet anfangen; darum ist auch Gottes Segen so reichlich bei ihnen. Wir sollten aber aus solcher Historien uns bessern und dergleichen und auch zu all unserm Thun gottseilig und christlich schicken. Aber jedermann thut das Widerspiel: Gott fürchten wir nicht, um seinen Segen bitten wir nicht, schelten, fluchen ist das erste; darum stehets und gehets so übel in aller Welt.

\section{Gebet}

Verleihe der treue GOtt, daß wir auch fleißig und herzlich beten und gottsfürchtig leben, und also aus Gnaden seinen Segen an Leib und Seele haben und behalten, bis zum freudenreichen ewigen seligen Leben, Amen.

\chapter{Das 25. Kapitel}

\section{Vorrede}

Mitten im Leben sind wir mitten im Tode. Denn wie wir von wegen der Sünden alle Augenblick mancherlei und tödtlichen Krankheiten

unterworfen sind: also ist kein Unfall so groß und schwer, er kann uns bald zu Haus und Hof kommen und uns unversehens dahin reißen. Darum bedürfen wir allezeit der Gnade Gottes, daß er uns behüte und erhalte. Das aber solches uns widerfahre, so sollen wir auch alle in christlicher Buße leben, auf daß uns der Tod in unsern Sünden nicht übereile, daß wir in Sünden zum ewigen Verderben nicht umkommen und sterben. Ferner höret an das 25. Kapitel.

\subsection{Abrahams andere Ehe. Ismaels und Isaaks Söhne}
\quellseite{80}

a Er fiel, mancherlei Deutung kann hie sein. Meine Deutung ist diese, daß Ismael ein herrlicher Mann gewesen sei, daß zu seinem Ende kommen und alle seine Brüder und Freunde, und ist für denselben ehrlich und löblich gestorben.

\section{Summarium}

Dies ist eine Historia, wie Abraham gestorben sei, v. 7.; und daß sein Sohn Isaak eben die Anfechtung hat müssen tragen, wie sein Vater, daß er zwanzig Jahr in der Ehe gewesen, und dennoch keinen Erben hat gehabt, darum, daß Rebekka unfruchtbar war, v. 21.

Weil nun aber die Verheißung auf Isaaks Erben stund, läßt er sich solches Aergerniß nicht abschrecken, sondern hält sich an Gottes Gebet, und erlanget einen, Zwillinge, Jakob und Esau, v. 26., welche gleich wie Ismael und Isaak bedeuten die zwo Kirchen: eine, so auf Gottes Barmherzigkeit traut; die andere, so auf eigen Werk und Frömmigkeit bauet.

Daß Esau die Herrlichkeit der Erstgeburt um eine Schüssel voll Linsen verkauft, v. 31., welche fürnehmlich hatte die Verheißung von Christo, ist ein sonderlich Exempel des Unglaubens, daß er solche Verheißung für nichts geachtet hat; Jakob aber, als ein gottseliger Mensch, achtet dieselbige hoch und theuer, darum erlanget er sie auch, und wird solches Glaubens halben Jakob vom Propheten Hosea Kap. 12, 3. herrlich gepriesen. Denn daß er in der Geburt den Esau bei den Fersen hält, v. 26., solches deutet Hoseas auf den Glauben Jakobs, daß er gern die Erstgeburt und die Verheißung vom Herrn Christo gehabt und seinen Bruder daran wollte gehindert haben.

\section{Votum}

Der barmherzige treue GOtt stärke unsern Glauben, daß wir sein heilig Wort lieben und üben, christlich leben, alles Creutz geduldig tragen und endlich selig sterben, hochgelobt in Ewigkeit, Amen.
\quellseite{81}

\chapter{Das 26. Capitel}

\section{Vorrede}

Der Welt Art und Gebrauch ist, daß sie die Christen verfolgt, beleidigt und beschädigt: Gottes Art und Weise dagegen ist, daß er die Christen vertheidigt, schützet und segnet. Derowegen wer den Welt-Kindern verhalten muß, der befehle es Gott, räche sich nicht selbst, bete und sei in seinem Beruf fleißig. So wahr als Gott im Himmel lebet, er wird dem Gerechten helfen und seine Stärke sein in der Noth; er wird die Gläubigen versorgen und wird sie nicht ewiglich in Unruhe lassen, Ps. 55. Dieses lehret auch das 26. Kapitel.

\section{Isaaks Wanderschaft, Wohnung und Widerwärtigkeit}

a Esek heißt Unrecht, wenn man jemand Gewalt und Unrecht thut.

b Sitna heißt Widerstand; daher der Teufel Satan heißt, ein Widerwärtiger.

c Rehoboth heißt Raum und Breite, das nicht enge ist.
\quellseite{82}

d Saba heißt ein Eid und Schwur, oder die Fülle, der aber heißt ein Brunn.

\section{Summarium}

Dies ist eine Historia, in welcher wir sehen, wie Isaak Theurung halber an einen fremden Ort ziehen, V. 1., und daselbst nicht allein seiner Rebekka halben in Gefahr leben, sondern auch allerlei Ueberlaß von den Philistern leiden muß, V. 7. Wir sollen daraus lernen, wie unser Gott seine Christen in der Welt unter bösen Buben in mancherlei Weise plagen läßt, aber doch ihnen auch väterlich hilft und sie schützet. Denn je mehr die Welt die Frommen will dämpfen, je mehr er sie segnet und ihnen hilft.

\section[Gebetum]{\unsicher{Gebetum}}

Gott steure und wehre den Verfolgern seiner Christenheit, und segne und behüte alle, welche die bedrängten Christen fördern, nähren und schützen, und thue solches um Jesu Christi seines lieben Sohns willen, Amen.

\chapter{Das 27. Kapitel}

\section{Vorrede}

Gott ist ein Brunnquell alles Guten denen, die ihn fürchten; dagegen ist er ein verzehrend Feuer allen, die ihn verachten. Derowegen wer von Gott will gesegnet sein, der diene ihm in wahrem Glauben, und dem Nächsten in wahrer Liebe, daß ist, er vertrauhe auf Gottes Verheißung und beschädige nicht seinen Nebenchristen; sondern rede und thue, das ihm in seinem Anliegen nöthig ist. Also wird er von Gott aus Gnaden bekommen und haben, was ihm und den Seinen gut und nützlich ist; und dagegen wird ihn Gott aus Barmherzigkeit behüten und schützen, wider alles, was ihm und den Seinen böse und schädlich ist. Das unterrichtet auch das 27. Kapitel.

\section[Jacob wird von seinem Vater gesegnet, von Esau]{Jacob wird von seinem Vater gesegnet, von Esau angefeindet}
\quellseite{83}

\section{Summarium}

Dies ist eine Historia, wie Esau zwei heidnische Weiber nimmt, ohne Zweifel wider seiner Eltern Rath und Willen, V. 1., und der Segen auf den frommen Jakob kommt, welchen Segen zuvor Esau aus dem Unglauben verachtet und um ein Essen Linsen ihm verkauft hatte, und zur Warnung; denn wer Gott und seinen Segen verachten kann, demselben soll der Fluch zu Theil werden.

Daß aber Jakob durch seine Mutter mit einer List zum Segen gefördert wird, V. 6., geschieht aus dem Befehl des Herrn, da er sagt: Der Größere wird dem Kleinern dienen.
\quellseite{84}
Maße sehr betrübt; und sprach zu seinem Vater: Segne *mich auch, mein Vater.

a) Untertreten) Ebr heißt eine Fußsohle. Daher kommt Jakob oder Jacob, ein Untertreter, oder der mit Füßen tritt, und bedeutet alle Gläubigen, die durch das Evangelium die Welt, das Fleisch und den Teufel mit Sünde und Tod unter sich treten durch Christum.

\section{Summarium}

Dies ist eine Historia, wie der Unglaube und die Verachtung der Verheißung Gottes am Esau leiblich gestraft wird. Denn auf dem erstgebornen Sohn stund das Regiment und die Verheißung von Christo. Weil nun Esau solches beides nicht größer achtet, denn ein Essen Linsen und Bauchfülle, nimmt Gott solchen Segen von ihm und gibt ihn den Jakob.

Dies ist nun die Ursache, daß Esau den Jakob verfolget. Wie in der Welt je und je geschehen ist, daß die rechte christliche Kirche, so sich Gottes Güte tröstet, von den Werk- und Bauchheiligen muß verfolget werden.

\section{Votum}

Der allmächtige treue GOtt steure allen Feinden seiner Christenheit, und erhalte seine Gläubige unter seinem Schutz, in Leibes und der Seelen Gefahr, bis zum ewigen Leben, Amen.

\chapter{Das 28. Kapitel}

\section{Vorrede}

Wenn die Kinder ihren treuen Eltern und die Eltern ihrem frommen Gott gehorsam sein, so gelinget es beiden. Denn die Kinder, so ihre Eltern ehren, die ehret Gott wieder an Leib und Seele, in diesem und künftigen Leben; also, die Eltern, die in Gottesfurcht leben, sich treulich nähren und Gottes Segen recht gebrauchen, die genießen es auch zeitlich und ewig, beides aus Gnaden, um Jesu Christi willen. Derowegen Kinder und Eltern fleißig auf Gott und seine Gebote sehen und denen treulich und beständig folgen sollen. Davon höret auch an das 28. Kapitel.

\section{Jakob ziehet auf seiner Reise die Himmelsleiter}
\quellseite{85}

a Deinen Samen) hie wird dem dritten Patriarchen Christus verheißen, der Heiland aller Welt, und das künftige Evangelium von Christo in allen Landen zu predigen durch die Engel auf der Leiter vorgebildet.

b Heilig) heilig heißt hie metuendus, terribilis, Nota, da man GOtt fürchten und ehren sollte, als er daselbst will gefürchtet und geehrt sein. Daher auch derselbe Berg, Morija, timor, reverentia, cultus Dei. heißt, sup. Kap. 22. 2. Denn Gottesfurcht ist der höchste Gottesdienst. Und ist hie angezeigt, wo Gottes Wort ist (wie Jakob hie hört) da ist Gottes Haus, da stehet der Himmel offen mit allen Gnaden.

c Mein Gott sein) nicht daß er vorhin nicht sein Gott gewesen sei, sondern er gelobet einen Gottesdienst aufzurichten, da man predigen und beten sollte; da will er den Zehnten zu geben den Predigern, wie Abraham dem Melchisedec den Zehnten gab.

\section{Summarium}

Dies ist ein fein Exempel, wie Jakob seinen Eltern folget und nach ihrem Willen sich begehret zu verheirathen, V. 1. Denn in solchem Falle die Kinder sonderlich auf der Eltern Willen sehen und sich befleißigen sollen, daß sie wider ihrer Eltern Willen und Vorwissen nichts fürnehmen. Wie man siehet, daß selten Glück dabei ist, wo die Kinder wider der Eltern Willen heirathen und die Eltern hoch betrüben. Denn Gott will Eheleuten Glück geben, wo sie Gott fürchten, Pf. 128, 1. 4. Wo man aber den Eltern nicht folget und wider ihren Willen thut, da kann keine Gottesfurcht sein, wird auch kein Glück da sein.

Darnach ist eine andere Historia, wie Gott dem Jakob erscheinet in dem Traum und bestätiget mit einer neuen Verheißung den Segen von Christo, und dem Lande Canaan, welches Esau verachtet hatte, V. 10. Darum verpflicht sich Jakob wieder gegen Gott, er wolle alle seine Hoffnung und Trost allein auf solchen gnädigen Gott setzen, ihm dienen und an dem Ort ein Gottes-
\quellseite{86}
haus bauen, und mit dem Zehenten aller seiner Güter den rechten Gottesdienst am selben Ort fördern und erhalten, V. 20. Und zum Exempel, daß wir auch lernen Gottes Segen also erkennen, denselbigen zu Erhaltung seines Worts und des rechten Gottesdienstes brauchen.

\section{Datum}

Helfe der treue barmherzige GOtt, daß alle Kinder und Eltern ihm allzeit gehorsam seyn, und segne beide an Leib und Seele, zu zeitlicher und ewiger Wohlfahrt, Amen.

\chapter{Das 29. Kapitel}

\section{Vorrede}

Wie man im Alten Testament alle Morgen frisch Holz zum Altar des HErrn bringen und damit das Feuer, daß es nimmermehr verlöschte, erhalten mußte: also sollen wir auch alle Morgen und Abend neu Holz vom Baum des Lebens, welcher ist unser HErr JEsus Christus, das ist, von seinem heiligen Wort, in unser Herz legen, damit das Feuer des Glaubens, der Liebe gegen Gott und unsern Nächsten, und des Gebets, welches der heilige Geist in uns anzündet, nicht ausgelöscht und vertilget werde. Wie denn ein Feuer auslöschen muß, wenn man kein Holz anleget. Denn wir sollen unser Licht brennen lassen, vermahnet Christus unser HErr, Luc. 12. Wir sollen uns auch unter einander brünstig lieb haben, lehret Petrus 1 Petr. 1. Und wir sollen mit Ernst Gott anrufen und mit dem Geist und Sinn beten, unterrichten David und Paulus, Ps. 145. 1 Cor. 14. Ferner höret an das 29. Kapitel.

\section{Jakob erwirbt durch seinen vierzehnjährigen Dienst zwei Weiber}
\quellseite{87}

\section{Summarium}

Dies ist eine Historia, wie Jakob auch in Ehestand kommen sei, V. 20. Auf daß man an der heiligen Patriarchen Exempel sehe, welche doch Gottes Wort geprediget, und mit Opfern als rechte Priester umgegangen und Gott gedienet haben, daß der Ehestand eine heilige Weihwohnung sei; welchen doch der Papst als unheilig seinen Geistlichen verboten hat. Zum andern ist ein fein Exempel einer sonderlichen Geduld und Einfalt des Jakobs: welcher, ob ihm wohl die andere Schwester ward zugelegt, solches Unrecht duldet, und noch sieben Jahr um die Rahel dienet, so ihm verlobet war, V. 28. Zum dritten stehet man hie, wie Kinder Gottes Gaben sind, welche Gott auch darum der Lea gibt, auf daß er sie ihrem Jakob lieb und werth mache, weil sie ungeschaffener war, denn die Rahel, V. 31. Denn Gott will, daß zwischen Eheleuten alle Liebe und Freundschaft gehen soll; gleich wie der böse Feind alles dahin wendet, daß er zwischen Eheleuten Uneinigkeit anrichten und die Herzen trennen möge.

\section{Votum}

Der gnädige gütige Gott segne alle christliche Eheleute an Kindern und Nahrung, verleihe auch ihnen Geduld unter ihrem Hauskreuz, und helfe, daß sie und alle ihre lieben Kinder in seiner Furcht leben und sterben, Amen.

\chapter{Das 30. Kapitel}

\section{Vorrede}

Das Alte Testament ist ein Gesetzbuch, das da lehret, was man thun und lassen soll, und daneben anzeiget Exempel und Geschichte, wie solche Gesetze gehalten oder übertreten sind. Gleich wie das Neue Testament ein Evangelium oder Gnadenbuch ist, und lehret, wo man es nehmen soll, auf daß das Gesetz erfüllet werde. Wie aber Moses nun rufet dem ganzen Israel und sprach zu ihnen: Höre Israel die Gebote und Rechte, die ich euch heute für euren Ohren rede, und lernet sie, und behaltet sie, daß ihr darnach thut, 5 Mos. 5.: also sagt unser Herr Christus zu allen seinen Christen: Sie haben Moses und die Propheten, laß sie dieselben hören, Luc. 16. Derowegen lasset euch die heilige Biblia, und also Gottes Wort, treulich befohlen sein und betet, daß ihr darnach züchtig, gerecht und gottselig allezeit lebet, Amen. Und darauf höret ferner an den ersten Theil des 30. Kapitels.

\section{Jakobs Kinder und Reichthum}
\quellseite{88}

a Naphtali heißt verwechselt, umgewandt, umgekehrt, wenn man das Widerspiel thut. Ps. 18, 27.: Mit den Verkehrten verkehrst du dich.

b Dudaim) fragst du selbst, was Dudaim sind: es sollen Lilien, es sollen Beeren sein, und niemand weiß, was es sein sollen. Es heißen etliche Judenkirschen, die in der Weizen-Ernte reif sind.

\section{Summarium}

Dies ist eine Historia, wie Gott Jakob gesegnet und ihm Kinder gegeben hat; daß aber die zwo Schwestern so peinlich nach Kindern thun, und eine die andere neidet, geschicht alles darum, daß sie es für einen sonderlichen Segen und Gnade von Gott achten, wenn Gott Eheleuten Kinder gibt. Sonderlich aber, weil dazumal der verheißene Same Christus zu warten war, durch welchen alle Völker sollten gesegnet werden. Solche Hoffnung hat ohne Zweifel die frommen gottsfürchtigen Matronen getrieben, daß sie groß Verlangen nach der Leibesfrucht gehabt, ob doch dermaleinst Gott die gnädige Stunde wollte geben, daß dem Teufel seine Gewalt genommen und allen Menschen von Sünden und Tod geholfen würde.

\section{Votum}

Der Sohn Gottes, unser Herr Jesus Christus, lasse ihm alle gottseligen Eheleute befohlen sein, und behüte sie für dem Teufel und für allem Herzleid, Noth und Schaden, Amen.

\section{Der andere Theil des 30. Kapitels}

\section{Vorrede}

Wer Glück und Gedeihen zu seiner Nahrung haben will, der trachte am ersten nach dem Reich Gottes und nach seiner Gerechtigkeit, der bete fleißig um Gottes Segen dazu, der habe auch das Vertrauen zu Gott, daß seine Arbeit nicht vergeblich sein soll, und der treibe daneben seinen
\quellseite{89}
Beruf ohne des Nächsten Betrug und Schaden. Denn fleißig arbeiten und sich treulich nähren, aber ohne Gottesfurcht, ohne Gebet, ohne wahren Glauben an Gott, und mit anderer Leute Nachtheil, das bringet der Nahrung endlichen Verderb und Untergang. Das ist gewiß und Amen. Ferner höret an den andern Theil des 30. Kapitels.

Meinen Fuß) Das ist, ich habe müssen laufen und rennen durch dünn und dick, daß du so reich würdest. Mein Fuß hat müssen thun.

Unde pedes Evangelizantium pacem et cursus verbi seu ministerii das ist, daher die Füße derer, die das Evangelium und den Frieden predigen und der Lauf des Wortes- oder Predigtamtes.

d Du mußt dich hie nicht irren lassen, daß Moses das kleine Vieh jetzt Ziegen, jetzt Lämmer, jetzt Böcke heißt, wie dieser Sprache Art ist. Denn er will so viel sagen, daß Jakob habe alles weiße Vieh behalten und alles bunte und schwarze Laban gethan. Was nun bunt von dem einfärbigen Vieh käme, das sollte sein Lohn sein. Es ward Laban froh, und hatte die Natur für sich, daß von einfärbigen nicht viel bunte natürlich kommen; aber Jakob half der Natur mit Kunst, daß die einfärbigen viel bunte trugen.

\section{Summarium}

Ist eine Historia, wie Jakob als ein Hausvater auch für sein Weib und Kinder sorget, und ein Gedinge mit Laban macht seines Lohnes halben, doch also, daß er alles stellet, nicht auf seine Arbeit, sondern auf Gottes Segen, und zum Exempel, daß wir also unsere Nahrung suchen sollen, daß weder Geiz noch Uebervortheilen bei uns sey, sondern neben unsrer treuen Arbeit alles auf Gottes Segen gestellet werde. Solche göttliche Fürsorge des heiligen Patriarchen Jakobs sollten die unachtsamen vertrumbischen Hausväter merken, welche nichts für sich sparen, alles anwenden und nach Weib und Kind, welche doch Gott ihnen zu nähren befohlen, auch Geschicklichkeit und andern Vortheil dazu verliehen, nichts fragen. Eine gräuliche Sünde ist solches und wird mit mancherlei Unrath gestraft, wie vor Augen.
\quellseite{90}
Zum andern ist auch ein fein Exempel hie, wie um eines frommen Knechts willen Gott alles Glück dem Herrn und ganzem Hause widerfahren läßt.

\section{Votum}

Regiere und führe der treue GOtt die Herrschaft und das Gesinde, daß ein jedes seines Berufs also abwarte, auf daß es GOtt zu Lob, ihrem Nächsten zu gut, und ihnen selbst zur Wohlfahrt diene und gelange, Amen.

\chapter{Das 31. Kapitel}

\section{Vorrede}

Wie wir uns gegen Gott halten, nach seinen Geboten: also hält sich Gott gegen uns mit seinen Verheißungen. Denn weil in allen Ständen, hohen und niedrigen wider seine Gebote sehr gesündigt wird, so entzieht er uns auch an vielen Orten seine Zusagungen, und also an Leib, Seele und Nahrung, Segen und Wohlfahrt. Derowegen wer da will, daß ihn Gott benedeie und behüte, wie er versprochen hat, der lebe und wandele auch in seinem Beruf, Amt, Handel, Ober und Handthierung, wie ihm Gott befohlen hat. Daher sagt der Prediger Salomon im 8. Kapitel: Wer das Gebot hält, der wird nichts Böses erfahren. Daneben höret an den ersten Theil des 31. Kapitels.

Jakob ziehet mit Weib und Kind von seinem Schwiegervater Laban.

* Stahl das Herz, Herz stehlen ist, betrüglich geredt, so viel, als etwas thun hinter eines andern Wissen.

\section{Summarium}

Ist ein schön Exempel, darin wir sehen, wie Gott alle Treue den Frommen reichlich belohnen, wiederum alle Untreue damit will strafen, daß er seinen Segen dazu geben will. Denn ob wohl Laban den frommen Jakob nicht mit Treuen mei-
\quellseite{91}
net, weil aber Jakob mit allen Treuen ihm dienet, läßt ihn Gott solcher Treu genießen, Laban muß aber seiner Untreue entgelten und soll ihm kein Vortheil helfen. Auf daß wir auch lernen treulich und ungefährlich handeln, so anders Gottes Segen soll bei uns sein. Es ist in aller Welt, an allen Orten, sehr viel armes Volk, das Nahrung halben weder hinter sich noch für sich kann. Wer solches Unfalls Ursache gewiß will wissen, mag es hie lernen: untreulich dienen sie, sind nachlässig, unachtsam, ungehorsam. Solche Sünde muß darnach an ihnen gestraft werden, wenn sie zu Haus ziehen und auch in den Stand gerathen, da sie Gesinde halten und mit denselben sich nähren müssen.

\section{Votum}

GOtt vergebe allen ihre Sünde, die Vergebung begehren und bitten; und verleihe, daß ein jedes willig und fleißig wandele in seinem Beruf, wie es GOtt in seinem Wort, treulich und ernstlich, lehret und erfordert, Amen.

\section{Der andere Theil des 31. Kapitels}

\section{Vorrede}

Treue in allerlei christlichen Händeln und Gewerben belohnet Gott mit Segen als ein barmherziger, gütiger Vater: Untreue strafet Gott mit Ungefallen als ein gerechter, wahrhaftiger Richter. Derowegen ein jedes in seinem Beruf und Nahrung treulich und ungefährlich handeln soll, so will ihn Gott auch segnen und für Gefahr behüten; oder will ihn unter dem Kreuz trösten und erhalten. Dies lehret der königliche Prophet David im 128. Psalm, und spricht: Wohl dem, der den Herrn fürchtet und auf seinen Wegen gehet. Du wirst dich nähren deiner Hände Arbeit, wohl dir, du hast es gut. Und St. Petrus in Geschichten der Apostel im 10. Kapitel sagt also davon: In allerlei Volk, wer Gott fürchtet und recht thut, der ist ihm angenehm. Zu diesem höret auch an den andern Theil des 31. Kapitels.
\quellseite{92}

\section{Summarium}

An ein schön Exempel, in welchem wir sehen, wo man anders fromm und treu ist, daß Gott die Seinen nicht lassen, sondern für Unglück bewahren will. Denn ehe Jakob weiß, daß ihm sein Schwäher Laban nacheilet, der Meinung, daß er ihm will Schaden thun, nimmt sich Gott des frommen Jakobs an, strafet und bedrohet Laban, daß er mit Jakob nicht anders denn freundlich handeln solle. Darum fleißige man sich nur mit allem Ernst, daß man fromm sey und mit gutem Gewissen mit den Leuten handle; so wird Gott sich dermaßen auch gegen uns erzeigen und uns schützen und retten, wie den gottseligen frommen Jakob.

\section{Gebet}

Wir danken Gott für allen Segen, den er von Anfang seinen Gläubigen gegeben hat; der segne und behüte uns alle ferner, doch nach seinem gnädigen väterlichen Willen und Wohlgefallen, Amen.

\chapter{Das 32. Capitel}

\section{Vorrede}

Freilich ists wahr und recht gesagt: Je größer Noth, je näher Gott. Denn solches lehret Gottes Wort, dieses bezeugen auch die Exempel der heiligen Gottes. Dergleichen erfahren wir solches alle Tage an den Gläubigen. Derowegen sei niemand unter seinem Kreuze ungeduldig, wie lange es auch währet und wie schwer solches ist; sondern ein jeder Kreuz-Bruder und jede Kreuz-Schwester bete und gebrauche gebührliche und christliche Mittel. Der

b Furcht) Jakob nennet die Gott Isaaks Furcht darum, daß Isaak gottfürchtig war und Gottes Diener.

c Gilead) heißt ein Zeugenhaufe, und bedeutet die Schrift, da viele Zeugnisse von Gott häufig innen sind.
\quellseite{93}
Gott, der andern bisher wunderbarlich geholfen, wird euch auch in seiner Noth \unsicher{verlassen}, denn er ist Gott über alle Noth. Und darauf höret an das 32. Kapitel.

\section{Jakobs Furcht, Gebet und Kampf}

\section{Summarium}

Das ist ein schön Exempel, wie Jakob sich für seinem Bruder Esau demüthiget, und weil er sich für ihn besorgen muß, seine brüderliche Mittel suchet, den Zorn damit abzulehnen, V. 3. Am allerersten hebt er die Sache mit Gott an, und danket ihm um seinen reichen Segen: Da ich, spricht er V. 10., über diesen Jordan ginge, hatte ich mehr nicht, denn diesen Stab, aber jetzt bin ich zwei Heere worden, das ist, ich habe einen großen Haufen Rind und Gesinde, und darnach einen großen Haufen Viehes, das alles dein Segen und Gabe ist und ich billig dir dafür danke. Darnach betet er, und erinnert Gott seiner Zusagung, daß er ihn gnädig für dem Zorn Esau behüten wolle, V. 11. Im letzten braucht er auch diese Mittel, daß er mit Geschenken des Bruders Zorn gedenkt zu versöhnen, V. 14. Also siehet man, wo die heiligen Leute ihrer Vernunft gebraucht haben, daß sie alles aufs einfältigste und zum Besten gethan haben.
\quellseite{94}

a Rang) im Hebräischen kommt ringen vom Staub her, als wenn zween mit einander ringen, daß der Staub sich erhebt, und dicke um sie wird, und lautet so viel: Es staubet ein Mann mit ihm, das ist, ein heftiger Kampf war es; das soll niemand verstehen, denn die Erfahrung.

b Israel) Israel kommt von Sara und heißt kämpfen oder überwältigen; daher auch Sar ein Fürst oder Herr, und Sara eine Fürstin oder Frau heißt, und Israel ein Fürst oder Kämpfer Gottes, das ist, der mit Gott ringet und angewinnet. Welches geschieht durch den Glauben, der so fest an Gottes Wort hält, bis er Gottes Zorn überwindet, und Gott zu eigen erlanget zum gnädigen Vater.

c Pniel) Pniel oder Pnuel heißt Gottes Angesicht oder Erkenntniß. Denn durch den Glauben im Streit des Kreuzes lernt man Gott recht erkennen und erfahren; so hats denn keine Noth mehr, so gehet die Sonne auf.

\section{Summarium}

Ist eine Historia, wie ein Engel mit Jakob ringt, und Jakob obsieget. Es zeuget aber Hoseas Kap. 12, 4. klar, daß dieser Engel nicht ein gemeiner Engel, sondern unser Herr Christus, der Sohn Gottes, selbst gewesen sei, und mit Jakob gerungen habe. Was aber solch Ringen sei gewesen, ist ungewiß, wiewohl aus dem Text abzunehmen ist, daß Jakob nicht allein mit der Leibes-Stärke, sondern viel mehr mit dem Glauben und Herzen gekämpft habe, und fest an der Verheißung Gottes gehalten, daß ihn Gott wieder heimbringen und seinem Samen nach ihm segnen wolle. Solcher Glaube hat dem Engel obgesieget, das ist, wie Jakob geglaubet, und der Zusagung Gottes sich getröstet hat, also ist es ihm auch widerfahren; Gott hats ihm nicht können abschlagen, er hat ihm müssen gewähren. Denn sein Wort und Zusagung ist ja, und kann nicht fehlen, wo man mit Herzen sich darauf verläßt.

\section{Votum}

Der treue barmherzige GOtt stärke unsern Glauben, und stehe uns bei, auf daß wir auch nach seiner Verheißung an Leib und Seele gesegnet, und zum ewigen Leben erhalten werden, Amen.

\chapter{Das 33. Kapitel}

\section{Vorrede}

Gegen Gott sollen wir uns als Kinder verhalten; und gegen unsern Nächsten sollen wir uns als Brüder und Schwestern erzeigen. Kinder aber ehren ihre Eltern; und einander dienen treue Brüder und Schwestern. Derowegen wollen wir Gott unsern Vater allezeit mit Mund, Herzen und guten Werken loben und preisen; und gegen unsern Nächsten, als unsern Brüdern und Schwestern, wollen wir uns allezeit freundlich und dienstlich, mit Rath und That, beweisen. Beides lehret der heilige Apostel Paulus. Denn 1 Cor. 10. sagt er: Was ihr thut, so thut es alles zu Gottes Ehre. Und an die Galater am 5. Kapitel spricht er: Durch die Liebe diene einer dem andern. Davon höret auch an das 33. Kapitel.

\section{Versöhnung Jacobs mit Esau}
\quellseite{95}

a Mählig) merkt daß Rechtgläubige und Werckheilige nicht können mit einander wandeln. Denn die Gläubigen fahren säuberlich mit stillem Geist, aber die Werckheiligen fahren stark mit Vermessenheit ihrer Werke in Gottes Gesetzen.

b Oder Schafe.

\section{Summarium}

Das ist ein fein Exempel, wie freundlich sich Jakob stellet, und demüthiget für seinem ältern Bruder Esau, V. 3. Und zur Lehre, daß wir nicht allein gegen leiblichen Brüdern, sondern auch ein jeder Christ mit dem andern aufs freundlichste leben sollen.

Daß aber Jakob einen Acker kauft im Lande Canaan, und da eine Hütte bauet und dem HErrn einen Altar zurichtet, V. 19., ist erstlich eine Anzeigung, wie er an der Hoffnung fest gehalten; und obgleich die Verheißung mit dem Lande Canaan sich lang verzogen hat, dennoch nicht hat gezweifelt, es werde seinem Samen zum Erbe werden. Darum, da er mehr nicht kann, kaufet er einen sondern Ort zum Gottesdienst, zum Zeichen, daß noch Gott sein Volk darein bringen und seinen Dienst darin anrichten werde. Darnach ists ein Exempel, wie die alten Väter Gottes Wort treulich geprediget und rechten Gottesdienst erhalten, und ohne Zweifel von den Heiden darum auch viel haben leiden müssen. Denn der Altar hat zu den zweien Stücken müssen dienen: daß erstlich daselbst Gottes Wort verkündiget und die Leute vom künftigen gebenedeyten Samen sind unterwiesen worden, und darnach die äußerlichen Gottesdienste mit dem Opfer am selben Ort sind ausgerichtet worden.

\section{Votum}

Verleihe der gütige und barmherzige GOtt, daß heute christliche Eltern auch sein Wort fördern helfen, und daß sie mit ihren Kindern allezeit in aller Gottseligkeit leben, und, wenn es GOtt gefällt, auch selig sterben, Amen.

\chapter{Das 34. Kapitel}

\section{Vorrede}

Also gehet es den Kindern Gottes auf Erden: wenn ein Kreuz bei ihnen aufhöret, so ist es ein Anfang eines andern und größern, und treibet immer eins das andere. Und widerfähret ihnen gleich, wie die Heiden schreiben von ihrem Hercule, welcher, wenn er dem gräulichen und schrecklichen Thiere, Hydra genannt, kaum ein Haupt hatte abgehauen, so wuchsen flugs drei andere Häupter an derselbigen Stelle: also ists ein unruhig und mühselig Ding um eines gottesfürchtigen Menschen Leben. Denn er ein stetiger Kreuzträger sein muß, und wenn er kaum eine Anfechtung überwunden hat, so ist bald da eine andere. Dessen ist ein sonderlich Exempel der heilige Patriarch Jakob, von welches großem Kreuze einem meldet das 34. Kapitel.
\quellseite{96}

\section{Die Schwächung der Dina verursachet ein Blutbad}
\quellseite{97}

\section{Summarium}

Hier ist eine Historia, wie der fromme Jakob an seinen Kindern viel Jammer und Herzeleid erlebt hat: daß ihm seine Tochter geschändet, V. 2., und die zween Söhne einen jämmerlichen Mord haben angerichtet, V. 25. Aber dies ist der Kreuze eines, das bisweilen Gott auch auf fromme Leute pflegt zu legen, da es doch weder an Zucht, noch guten Exempeln mangelt, wie man in der Historie Davids wohl siehet.

Neben dem siehet man, wie Gott der Hurerei so feind ist, weil er eine ganze Stadt um solcher Sünden willen erwürgen läßt. Derohalben denn weltliche Obrigkeit mit allem Ernst solchem Laster wehren und in dem Fall Gottes Zorn und Strafe, der es nicht will ungestraft lassen, zuvor soll kommen. Aber leider, es gehet solche Sünde je länger je gräulicher, und schier ohne alle Scheu und Schande; darum muß es allenthalben desto übler und schwerlicher zugehen.

\section{Votum}

GOtt steure und wehre allen Lastern; und regiere die liebe Jugend, daß sie zu seinem Ehren, und ihren Eltern zur Freude und Trost, allzeit christlich und züchtig leben, Amen.

\chapter{Das 35. Kapitel}

\section{Vorrede}

Es ist heuer nicht neu, daß einem Christen oft alle Wochen, ja alle Tage ein neu Kreuz zu Hause und Hof kömmt. Denn wie man im Sprüchwort sagt: Je lieber Kind, je größer Ruthe; also: Je frommer Christ, je mehr, größer, länger und schwerer Kreuz er tragen muß. Davon redet der heilige Paulus also in Geschichten der Apostel im 14. Kapitel: Durch viel Trübsal muß man in das Reich Gottes gehen. Und an die Ebräer am 12. Kapitel wird gesagt: Welchen der HErr lieb hat, den züchtiget er; er stäupet aber einen jeglichen Sohn, den er aufnimmt. Wohl aber denen, die ihr Kreuz in Geduld tragen. Denn wer sein Kreuz trägt in gläubiger Geduld, der lebt und bleibt in Gottes väterlicher Huld. Davon höret auch an das 35. Kapitel.

\section{Rahel stirbt an der Geburt}

a Ohren-Spangen) Lunulas heißt man güldene Haar-Bänder. Nun sind's Perlen-Borten geworden, Sprüchw. 11, 23.
\quellseite{98}

b Trankopfer) das war Wein; wie das in den folgenden Büchern genugsam gesehen wird.

c Benoni) heißt meines Schmerzens Sohn.

d Benjamin) heißt der Rechten Sohn.

\section{Summarium}

Ist eine Historia, wie Jakob ein recht bischöflich Amt führet, daß er die Abgötterei unter den Seinen ausrottet, und nach Gottes Befehl den rechten Gottesdienst anrichtet, B. 2. Zum Exempel aller Obrigkeit, daß sie wisse, es sei ihr fürnehmstes Amt und höchstes Werk, über dem rechten Gottesdienst halten, und Abgötterei mit Gewalt ausrotten.

Zum andern siehet man, wie der fromme Vater Jakob wiederum zwei neue und schwere Kreuze erfahren muß: daß ihm seine liebe Rahel in Kindesnöthen stirbt, B. 18., und sein ältester Sohn Ruben sein Kebsweib beschläft, B. 22. Denn Gott will seine Heiligen auf mancherlei Weise hier auf Erden üben und versuchen. Solche Exempel mögen wir uns zum Trost merken, wenn dergleichen uns begegnet, daß wir gedenken, wir sind nicht allein, es sei dergleichen auch den großen trefflichen Heiligen widerfahren.

\section{Votum}

Der treue barmherzige GOtt lasse ihm alle christliche Kreuzträger zu Trost und Hülfe befohlen sein, und behüte uns alle für allem Unfall Leibes und der Seele, Amen.

\chapter{Das 36. Kapitel}

\section{Vorrede}

Dies Kapitel meldet nur, wie Esau gewachsen und zum großen Volke worden, und auf Erden mit Gewalt und Herrlichkeit sehr habe zugenommen. Ist derohalben unnöthig, daß man die bloßen Namen in der Kirche lese; es geht aber dem Sprüchwort nach: Unkraut wächst mit Haufen, und verdirbt nicht; da es dagegen mit dem frommen Jakob, der doch den Segen und die Verheißung hatte, sehr geringlich sich anließe und schwächlich fort ginge; Ursache der Satan, der ist den Frommen Feind und hindert, wo er kann.

\section{Geschlechtsregister Esaus}
% Scan 99: no retained apparatus prose; Bible text and other excluded blocks are inventoried in selection.json.
\quellseite{100}
hieß Mehetabel, eine Tochter Matreds, die Mesahabs Tochter war.

\chapter{Das 37. Kapitel}

\section{Vorrede}

Das 36. Kapitel des ersten Buches Mose lehret, wie Esau gewachsen und zum großen Volke worden, und auf Erden mit Gewalt und Herrlichkeit habe zugenommen (denn Unkraut wächst mit Haufen und verdirbet nicht); da es dagegen mit dem frommen Jacob, der doch den Segen und die Verheißung hatte, sehr geringlich sich anließe und schwächlich fortginge. Denn der Satan ist den Christen feind und hindert, wo er kann. Das lässet aber Gott darum geschehen, auf daß der Glaube geübet und geprüft werde, daß wir lernen allein am Wort hangen, und was sichtbar und greiflich ist, nicht ansehen, sondern uns daran halten, daß man nicht siehet, und daß wir den Trost behalten, davon St. Paulus sagt 2 Cor. 4.: Wir leiden Verfolgung, aber wir werden nicht verlassen. Wir werden unterdrückt, aber wir kommen nicht um. Darauf höret ferner an das 37. Capitel.

\section{Joseph wird aus Neid von seinen Brüdern verkauft}
\quellseite{101}

a Vater) das war Jaakob, Jakobs Vater und Josephs Großvater, welcher nach dieser Geschichte noch zwölf Jahre gelebet hat.

\section{Summarium}

Hie hebt sich an die Historia vom Joseph, in welcher wir sehen, wie dem frommen Jakob seine Söhne ein neues Herzeleid zurichten, daß sie Joseph verkaufen und sprechen, es habe ihn ein wildes Thier zerrissen, V. 20.

Dies soll uns ein sonderlich Exempel sein, an welchem wir lernen und sehen sollen, wie Gottes Rath durch Menschen Anschläge nicht möge gehindert werden. Joseph wird darum von seinen Brüdern verkauft, auf daß sein Traum nicht erfüllet werde, und sie ihn nicht müssen als einen Herrn anbeten. Aber eben mit solchem Rath fördern sie Joseph zu solcher Herrlichkeit, zu welcher er nimmermehr hätte mögen kommen, wo er bei dem Vater wäre blieben. Derohalben sollen wir nicht so bald verzagen, wenn das Kreuz kommt und es übel zugehet. Man harre nur; es kann Gott treffliche und große Dinge dadurch ausrichten, so man in der Furcht Gottes bleibet (und durch die Sünde nicht Ursache zur Strafe gibt), wie Joseph thut.

\section[Notum]{\unsicher{Notum}}

Der treue Gott verleihe solches heute über allen christlichen Kreuzträgern; und segne und behüte unsere liebe einheimische und ausländische Jugend, an Leib und Seele, um JEsu Christi, seines geliebten Sohnes willen, Amen.

\chapter{Das 38. Kapitel}

\section{Vorrede}

Der heilige Apostel Paulus warnet und spricht 1 Cor. 10.: Wer sich lässet dünken, er stehe, mag wohl zusehen, daß er nicht falle. Hat diese Warnung jemals vonnöthen gethan, so thut sie ja zu dieser gegenwärtigen Zeit vonnöthen, da die Aergernisse so gar groß und die Teufel so er
\quellseite{102}
schrecklich grimmig sind. Und: haben große hochbegabte Leute gestolpert, und schwere Sündfälle gethan, so mögen sich wohl vorsehen andere, die weniger und geringere Gaben haben, daß sie nicht auch in grobe und schwere Sünde gerathen. Derowegen, ach laßt uns wachen und beten, daß wir nicht in Anfechtung und Sünde wider das Gewissen fallen. Darauf höret weiter an das 38. Kapitel.

\section{Juda Blutsschande mit der Thamar}

a Hirten Mag auch heißen Freund, darnach die Puncta im Hebräischen sich setzen lassen. Denn Juda hat ja müssen Weide haben, vielleicht auch eines Freundes nicht gerathen mögen. Puncta können so wohl fehlen als treffen.

b Hirten oder Freund, welches du willst.
\quellseite{103}

\section{60	Das 1. Buch Mose 38. 39}

\section{Summarium}

Es ist eine Historia, welche am allermeisten von Mose geschrieben ist um des HErrn Christi willen, welches Geschlecht von Juda und Thamar herkommt. Nun war Thamar eine Heidin, und ist doch unsers HErrn Christi Großmutter, anzuzeigen, daß auch die Heiden Gottes Volk sollen werden. Ja Gott ist nicht allein der Juden, sondern auch der Heiden Gott. Nun werden aber solche Historien von den heiligen Vätern nicht darum geschrieben, daß man ihren Sünden folgen soll, sondern daß wir an ihnen erkennen sollen unsere verderbte Natur, und derhalben nicht sicher sein, sondern beten, und in Gottesfurcht wandeln.

Darnach, daß man lerne wider das böse Gewissen sich mit Gottes Barmherzigkeit trösten; sintemal Gott solche Sünder nicht hat weggeworfen, sondern in seiner Linie behalten. Aber hie gehört auch zu, daß man sich bessere und von Sünden ablasse, wie Juda thut. Denn seine zween Söhne Er und Onan, da sie sich nicht wollen bessern, tödtet sie der HErr ihrer Sünden halber, V. 7. 10.

\section{Votum}

Der barmherzige GOtt regiere uns, damit wir nicht Aergerniß anrichten, und helfe, daß wir uns seiner väterlichen Barmherzigkeit in allen Nöthen freuen und trösten, und alle selig werden, Amen.

\chapter{Das 39. Capitel}

\section{Vorrede}

Wie zweierlei Christen sind, was den Glauben anlangt, nämlich, starke und schwache: also sind auch zweierlei Christen, was das Leben und die Werke betrifft. Etliche sind stark, etliche sind schwach. Starke sind, welche sich im Leben also verhalten, daß sie mit äußerlichen Sünden und Gebrechen, die man augenscheinlich sehen kann, sich nicht beflecken. Schwache sind, die bisweilen, nicht aus bösem Vorsatz, sondern aus angeborner menschlicher Blödigkeit, und von ihrem eigenen Fleisch und Blut übereilet, fallen in Gebrechen, die sie denn auch alsobald wiederum erkennen, herzlich Reue und Leid darüber haben, Gott im Namen Jesu Christi um Verzeihung bitten, auch, mit Hülfe und Beistand des heiligen Geistes, forthin ihren Wandel vorsichtiger anzustellen, damit sie niemand muthwilliger Weise ärgerlich sein, sich befleißigen. So wolle Gott die Starken in ihrer Frömmigkeit bestätigen und erhalten, und die Schwachen stärken, auf daß sie erfüllet werden mit Früchten der Gerechtigkeit, die durch Jesum Christum geschehen, zur Ehre und Lobe Gottes, Amen. Darauf höret nun an das 39. Kapitel.

\section{Josephs Dienst, Keuschheit und Gefängniß}
\quellseite{104}

a In dem du) Sonst mußte er auch die Frau versorgen, eben so wohl als das ganze Haus, ohne daß er nicht bei ihr schlafen sollt.

\section{Summarium}

Es ist eine sehr schöne Historia, in welcher wir am ersten lernen, wie Gott fromme Leute segnen und Glück und Heil um ihretwillen auch andern geben will, V. 3. 5. Darnach wie Joseph sich keusch und rein hält, da ihm doch Ursach, Uebels zu thun, auf das nächste bracht ward, V. 8. Solch Exempel, daß er Gott für Augen hat und für Sünde sich hüte, ist der Anfang alles Segens und Gnade, so Gott dem Joseph hat widerfahren lassen. Denn wie im ersten Gebot stehet, will Gott die, so ihn lieben, das ist, so sein Wort für Augen haben, in tausend Glied segnen. Aber man lerne hie auch die Weise, so Gott mit seinen Heiligen hält, welche er will segnen, und etwas Sonderes aus ihnen machen. Denn er steckt sie zuvor in Leiden und Unglück bis über die Ohren, und ob sie gleich unschuldig sind, stellet er sich, als kenne er sie nicht, und frage nicht darnach, wie man hie an Joseph siehet, der fromm und unschuldig ist, und dennoch muß er ins Gefängniß, daß böse unzüchtige Weib bleibet in Ehren und Würden. Aber man halte nur fest über Gottes Wort und Willen und leide geduldig. Denn gewiß ist es, ob man schon über einer rechten Sache gestöcket und geplöcket wird, so soll doch die Wahrheit endlich an Tag kommen, und Gott will endlich dem gerechten Theil helfen, wo man anders seine Hülfe durch Geduld erharren will. Darum gibt Gott auch Joseph im Gefängniß Glück und Heil.

\section{Votum}

Gott der Vater aller Barmherzigkeit lasse ihm uns alle befohlen sein; und helfe, daß die Jugend mit uns, in Lieb und Leid christlich lebe, und zeitlichen und ewigen Segen, auch mit uns habe und behalte, Amen.

\chapter{Das 40. Kapitel}

\section{Vorrede}

Der König und Prophet David sagt im 34. Psalm: Der Gerechte, das ist, der gläubige Mensch, muß viel leiden. Und St. Paulus spricht in den Geschichten der Apostel im 14. Kapitel: Durch viel Trübsal müssen wir ins Reich Gottes eingehen. Wenn man nun glauben könnte, daß Gott den Menschen, die er am liebsten hat, das größte, schwerste und langwierigste Kreuz zuschickt, so würde viel Klage und Ungeduld unter dem Kreuz dahinten bleiben. Derowegen wer heute ein Kreuzträger ist, der gedenke, wie jener erfahrene, versuchte und christliche Kriegsmann recht und wohl gesagt
\quellseite{105}
hat: Es muß niemand mehr Gefahr und Unglück erwarten und aufsehen, als der, welcher Gott und dem Kaiser treulich dienet. Er freue sich auch, daß Gott unter dem Kreuz trösten und Geduld verleihen wolle, und daß auf das Kreuz endlich Hülfe und ewige Freude folgen werde. Dies erinnert auch das 40. Kapitel.

\section{Joseph leget den Gefangenen ihre Träume aus}

\section{Summarium}

Sie ist ein fein Exempel, welches dazu dienet, daß man sich in Gottes wunderbare Werke recht lerne schicken, V. 3. Denn Joseph muß in den Kerker und ist doch unschuldig. Solches wird ihm ohne Zweifel haben wehe gethan; aber weil er solch Unrecht geduldig leidet, und sein Vertrauen auf Gott setzet, er werde letztlich helfen und retten, muß solcher Kerker und Gefängniß ihm dazu dienen, daß man Gottes Gabe an ihm lerne erkennen, um welcher willen er hernach zu großen Ehren kömmt. Solcher Gabe braucht er also, daß er doch Gott allein die Ehre gibt.

Aber er muß hie zum andernmal die Untreue der Welt erkennen lernen, daß der Schenk sein vergißt und er noch zwei Jahr muß im Gefängniß bleiben, V. 23. Auf daß wir lernen fest auf Gottes Hülfe trauen, ob er sie gleich lange verzeucht.
\quellseite{106}

\section{Votum}

Der GOtt alles Trostes tröste alle die das Kreuz tragen, und behüte uns für Gefängniß, und für zeitlicher und ewiger Noth, um JEsu Christi, seines lieben Sohns und unsers HErrn willen, Amen.

\chapter{Das 41. Kapitel}

\section{Vorrede}

Wer droben im Himmel mit den Heiligen und Auserwählten in Freuden leben will, derselbe Mensch soll allhie auf Erden, nach dem Exempel der Heiligen und Auserwählten, in Gottesfurcht leben. Denn, die dem Teufel zu Willen leben in dieser Welt zeitlich, die werden mit ihrem Unwillen, beim Teufel auch in der Hölle, leben und leiden ewiglich. Derowegen fürchtet den HErrn und dienet ihm treulich und rechtschaffen. Und so euch der Satanas euer Widersacher mit Versuchung und Anfechtung zusetzt, so thut ja nicht seinen Willen, sondern widerstehet ihm feste im Glauben, und leget nimmer ab den Harnisch Gottes, das ist, lebet in täglicher Buße, daß ihr bestehen könnt gegen die listigen Anläufe des Teufels und das Feld behalten möget. Dies betrachten wir heute auch bei dem 41. Kapitel.

Joseph wird durch Auslegung der Träume Pharaos zu fürstlicher Hoheit erhoben.

a Bei mir nicht) will sagen: ich bins nicht, der die Träume gedeutet hat oder könne; Gott ists, der es durch mich gethan hat, kann dirs auch thun.
\quellseite{107}

\section{Summarium}

Es ist eine Historia, in welcher wir an Joseph lernen: ob Gott schon die Frommen eine Zeit lang läßt leiden, daß er sie doch endlich errettet, und ihnen das Unglück zum besten gedeihen muß. Denn wo Joseph nicht wäre im Kerker gewesen, wäre mans nicht von ihm innen worden, daß er solche Gabe hätte, die Träume zu deuten, wäre auch nicht zu solchen Ehren kommen. Darum sollen wir lernen, daß wir uns am Kreuz nicht ärgern, sondern auf Gottes Hülfe sehen, und derselben gewiß warten. Neben diesem ist auch eine nützliche Hauslehre hie, V. 47., daß die Hausväter, wenn Gott Segen zur Nahrung gibt, solchen Segen nicht verschwenden sollen, wie doch gemeiniglich geschieht. \unsicher{Steermann} genennt, weil man viel hat, man soll viel umbringen und aufheben lassen. Aber Gottes Segen sollen wir aufheben und auf künftige Noth für uns und auf andere sparen. Denn weil der meiste Theil Gottes Segen mißbrauchet und undankbar ist, folget gemeiniglich auf eine große Wohlfahrt große Theurung, da soll ein vernünftiger Haushalter, wie Joseph denn hie that, sich aufrüsten.

b) Abrech) was Aelters heiß, lassen wir die Zänker \unsicher{suchen} an den jüngsten Tag, wollen \unsicher{vielmehr} verstehen wie es verdeutscht ist.
\quellseite{108}

\section{Summarium}

Ist ersichtlich ein fein Exempel, daß wir, wie Joseph, unsere Gaben nicht sollen brauchen zu eigener Hoffahrt und Nutz, sondern zu des Nächsten Wohlfahrt. Denn darum hat Gott solchen Traum Pharao sehen und die Deutung Joseph sehen lassen, weil eine solche Hungersnoth fürhanden war, daß die Leute sich darein schicken und der guten Jahre desto besser zum Vorrath brauchen sollen.

Darnach ist ein fein Exempel einer feinen fürstlichen Tugend, daß Pharao den hochbegnadeten Mann Joseph theuer achtet, ihm einen Vater des Landes nennt und den nächsten nach dem Könige ordnet, B. 43. Denn in Regimenten soll man sonderlich darauf Achtung geben, daß tapfere, geschickte Leute in Ehren gehalten werden. Wo aber leichtfertiges, liederliches Gesinde bei den großen Herrn lieb und werth ist, da hauset man selten wohl.

\section{Votum}

Verleihe GOtt, daß Obrigkeit und Unterthanen nach seinem Willen leben, und daß ein jedes endlich von allem Uebel erlöset, mit ihm in ewiger Sicherheit und Freude lebe, Amen.

\chapter{Das 42. Kapitel}

\section{Vorrede}

- Auf Gottesfurcht folget ein fröhlich Gewissen und Wohlfahrt; auf die Sünde folget dagegen ein böses Gewissen und Unglück oder Strafe. Denn von der Furcht Gottes saget also Sirach im 1. Capitel: Die Furcht des HErrn macht das Herz fröhlich und gibt Freude und Wonne ewiglich. Wer den HErrn fürchtet, der wird endlich den Segen behalten. Von der Sünde aber spricht David im 32. Psalm: Der Gottlose hat viel Plage; und Esaias im 57. Kapitel: Die Gottlosen haben nicht Friede, spricht mein Gott; derowegen laßt uns in der Furcht des HErrn leben und die Sünde fliehen, daß wir nicht mit der äußerlichen Züchtigung Seelen-Angst und des Gewissens Plagen und Nag[e]n ewig müssen tragen. Ferner höret an das 42. Kapitel.

\section{Reise der Söhne Jakobs in Egypten, ohne Benjamin}
\quellseite{109}

\section{Summarium}

Es ist eine feine Historia, wie Gott seine Räthe heimlich und wunderlich führet, bis er es dahin bringet, darauf sein Wort verrichtet hat. Denn hie sehet man, wie wunderbarer es sich spielet, ehe denn der Traum erfüllet wird, welchen Joseph seinen Brüdern gesagt, und sie ihn gern gehindert und ihn darum verkauft hatten. Aber eben damit sie es gedenken zu hindern, fördern sie es. Denn Gottes Gedanken müssen für sich gehen; der Menschen Gedanken aber müssen hinter sich gehen.
\quellseite{110}

\section{Summarium}

Es ist ein fein Exempel an den Brüdern Joseph, wie sie ihr Gewissen naget und sie jetzt an ihre Sünde gedenken, da sie es längst sollten vergessen haben. Aber die Sünde läßt sich nicht vergessen, sondern muß das Herz unruhig machen, es geschehe über lang oder über kurz. Darum mag man sich dafür hüten. Es war bisher unter den Brüdern Joseph alles still gewesen; aber jetzt mußte Simeon gefangen liegen und mit den andern das Gewissen und Sorge tragen, Gott wollte das Blut Josephs von ihnen fordern. Das Geld in Säcken, V. 35, macht ihnen auch Beisorge; ihre Untugend mußte gestraft werden. Endlich beklagt sich auch Jakob, ihr Vater, daß sie ihn um seine Söhne bringen, V. 36. Also wird ihnen mancherlei Weise ihre Sünde ins Gewissen eingetrieben, alles uns zur Warnung, daß wir’s bedenken und nicht sündigen, noch in Sünden sicher sein.

\section{Votum}

Gott der heilige Geist regiere uns, daß wir Gottes Wort folgen, und für Sünden und hüten, auf daß wir nicht Sündenstrafen leiden dürfen, sondern in Gottes Gnade leben und sterben, Amen.

\chapter{Das 43. Kapitel}

\section{Vorrede}

Gottes Wort wird täglich gelehret; aber viel hören es also, daß sie sich dennoch mit Sünden in ihrem Beruf, Handel, Erwerb und Haushaltung nähren, und daß sie unter dem Kreuz wenig zu ihrem Trost an Gottes Verheißung gedenken. Derowegen sehet zu, wie ihr Gottes Wort höret, und seid nicht allein Hörer, sondern auch Thäter desselben, und tröstet euch auch damit. Denn bei denen, die den Bund halten und darnach thun, bleibet Gottes Gnade ewiglich, sagt David im 103. Psalm. Und St. Paulus Röm. 2. spricht: Für Gott sind nicht gerecht, die das Gesetz hören, sondern die das Gesetz thun. Und Luc. 12. sagt auch Christus: Der Knecht, der seines Herrn Willen weiß und hat nicht nach seinem Willen gethan, der wird viel Streiche leiden müssen. Und Ebr. 4. wird gesagt: Das Wort der Predigt half jene nichts, da nicht glaubten die, so es höreten. Dazu höret an das 43. Kapitel.

\section{Reise der Söhne Jacobs in Egypten mit Benjamin}
\quellseite{111}

\section{Summarium}

Es ist zum ersten ein Exempel, wie auch heilige Leute, die die gemeinen Strafen der Welt, als Hunger, Krieg und andere dergleichen, müssen mit tragen, und es GOtt dennoch also mit ihnen schaffet, daß sie sonderlich Hülfe und Schutz finden, B. 1.

Darnach ist’s auch ein fein Exempel, wie Jakob treulich und ehrbarlich handelt, daß er nicht auf seinem Nutz, sondern auf die Billigkeit siehet, und das Geld so in Säcken gefunden ward, heißt wiederbringen, B. 12. So man solches nach der jetzigen Welt Art wollte rechnen, würde man wenig Leute finden, die es ihm nachthäten; sogar ist jetzt alles im Geiz und Eigennutz ersoffen.
\quellseite{112}

\section{Summarium}

Dies ist ein Exempel, wie Gottes Zusagungen müssen erfüllet werden und unverhindert bleiben, es thue die Welt dazu, was sie wolle. Denn hie wird Josephs Traum recht erfüllet, von den Garben und Sternen, weil seine Brüder für ihm sich neigen und niederfallen, und hilft sie nichts, was sie, solchem vorzukommen, fürgenommen hatten.

Darnach ist auch wohl zu merken die Freundlichkeit, so Joseph gegen seine Brüder brauchet, B. 27. Denn auch in heiligen Leuten die Natur ihre angeborne Eigenschaft behält, daß Brüder sollen gegen einander freundlich sein. Darum ist es eine unnatürliche, teuflische Untugend, wo Brüder untereinander uneins oder unfreundlich sind.

\section[Notum]{\unsicher{Notum}}

Helfe der treue GOtt, daß wir alle einander als Brüder und Schwestern herzlich lieben, damit er aus Gnaden uns wieder liebe, und die Werke seiner Liebe erzeige, um JEsu Christi, seines geliebten Sohns, unsers HErrn willen, Amen.

\chapter{Das 44. Kapitel}

\section{Vorrede}

Gott züchtiget und strafet uns, nicht, daß er uns verderbe, sondern, daß er uns dadurch zur Buße berufe, auf daß wir also gebracht werden zum Erkenntniß der begangenen Sünden, und daß dadurch den künftigen Sünden gewehret, und daß wir daneben zur Gottseligkeit gebracht und geführt werden. Davon sagt das Buch der Weisheit im 2. Kapitel also: Da strafest säuberlich, so da fallen, und erinnerst sie mit Zucht, wenn sie sündigen, auf daß sie von der Bosheit los werden und an dich, HErr, glauben. So pflegt man auch sonst zu sagen: Das Kreuz entweder den Menschen probirt, oder demütiget, oder reiniget, oder Gottes Ehre verkläret. Derowegen ein jeder sein Creuz in Geduld tragen und zur Buße und Uebung des Glaubens gebrauchen soll. Davon lehret auch das 44. Kapitel.

\section{Josephs Brüder werden hart geängstiget}
\quellseite{113}

\section{Summarium}

Sie ist eine sehr feine Historia, wie Joseph seinen Brüdern, ehe er sich zu erkennen gibt, einen neuen Schrecken einjaget mit dem gestohlenen Becher, B. 2. Er thut aber solches der Meinung, auf daß er ihnen die Sünde wohl eintreibe und das Gewissen weidlich rühre, daß sie aus Neid ihn erstlich hätten tödten wollen, und darnach unter die Fremden verkauft, und ihren frommen Vater so schmerzlich betrübet hätten. Solche Sünde hatte ihr Gewissen nicht dermaßen daheim gesucht, wie jetzt geschieht.

\section{Votum}

Der treue Gott erleuchte uns durch den heiligen Geist, daß wir unsere Sünde erkennen, die Gott abbitten, fliehen und meiden, damit wir in wahrer Buße christlich leben und selig sterben, Amen.
\quellseite{114}

\chapter{Das 45. Kapitel}

\section{Vorrede}

Das fünf und vierzigste Kapitel des ersten Buchs Mose ist ein recht tröstlich Kapitel und möchte wohl ein Brunnquell reichen Trostes genannt werden. So helfe der treue, liebe Gott, daß dadurch erquicket werden aller derer Herzen, die auch ihrer Sünden halben Leid tragen, und derer, die auch Kreuzträger und Kreuzträgerin sind. Denn wahrhaftig, Gott tröstet unter dem Kreuz, und hilft endlich von allem Kreuz. Und Gott versucht also seine Gläubigen und schicket ihnen ja zu Anfechtung und allerlei Noth und Widrigkeit, aber doch also, daß er sie wiederum erquicket und aufrichtet mit großem, starkem Trost, und daß er sie erfreuet mit nötiger und großer Hülfe. Davon höret auch an das 45. Kapitel.

\section{Joseph gibt sich seinen Brüdern zu erkennen}

a Lasset euch euren Hausrath nicht hindern, was ihr nicht verkaufen könnt in solcher theuren Zeit, das laßt hinter euch.
\quellseite{115}

\section{72 Das 1. Buch Mose 45. 46}

† Das hebräische Wort heißt eigentlich anders thun und anders werden, Thren. 2, 18. Fleto et non desinat pupilla oculi tui, neque quiescat pupilla oculi tui. Weine, und laß deine Augen nicht anders thun, Kap. 3, 49. Ich kann nicht anders thun. Lex Tapug. Hab. 1, 4. Es gehet anders denn recht, Recht gehet anders, gilt nichts. Also auch Jakob gedenkt viel anders, als sie es ihm erzählen. Sic Jacob longe aliud sentit, quam illi narrant.

\section{Summarium}

Es ist eine sehr feine Historia, wie Joseph sich seinen Brüdern zu erkennen gibt, B. 3, und thut wie ein rechtschaffener Prediger thun soll. Nachdem seine Brüder bekümmert und ihrer Sünden halben erschrocken waren, daß er sie wieder tröstet und ihre Sünden, so sie an ihm begangen hatten, zum Besten deutet. Er schilt sie nicht, gedenkt auch ihres Neides und Hasses nicht. Denn solches alles hatte zuvor ihr eigen Gewissen schon gethan, darum war es Zeit Tröstens. Sondern zeiget ihnen aufs allerfreundlichste an, wie Gott durch ihr Verkaufen ihn dahin gefordert habe, daß er sie und ganz Egypten von der Hungersnoth erretten sollte, B. 6. Aber solches hat Joseph selbst hernach müssen lernen; da er verkauft und hernach lange Zeit gefangen war, wußte er’s noch nicht. Auf daß wir auch lernen des Guten hoffen, und nicht kleinmüthig seien, wenn es uns übel gehet. Denn wir wissen nicht, was Gott damit meinet, und daß es uns zum Besten gedeyen soll, wie Paulus sagt Röm. 8, 28.: Den Gottseligen soll es alles zum Besten dienen.

\section{Votum}

Gelobet und gepreiset sei der gütige, treue GOtt, daß er betrübt und tröstet, daß er Kreuz aufleget, auch Kreuz tragen hilft und wegnimmt; der wolle uns auch in aller Noth allzeit trösten und helfen, Amen.

\chapter{Das 46. Kapitel}

\section{Vorrede}

Von Gottes Verheißungen fällt nicht ein einzig Wort auf die Erde, sondern sie werden alle endlich und vollkommenlich erfüllet. Denn also sagt St. Paulus 2 Cor. 1.: Alle Gottes-Verheißungen sind ja in Christo, und sind Amen in Christo, Gott zu Lobe. Derwegen laßt uns nach Gottes Wort und Geboten gehorsamlich leben und wandeln, wie es eines jeden Beruf und Stand erheischt; so wird sich Gott nach seinen Zusagungen und Verheißungen gegen uns väterlich verhalten und erzeigen, wie es unsre Nothdurft und Wohlfahrt in diesem und künftigem Leben erfordert. Denn die Gottseligkeit ist zu allen Dingen gut und hat die Verheißung dieses und des zukünftigen Lebens, sagt St. Paulus 1 Tim. 4. Darauf höret auch an das 46. Kapitel.

\section{Jakobs Reise nach Egypten zu seinem Sohne Joseph}
\quellseite{116}

\section{Summarium}

Hier ist eine Historia, wie Jakob sei in Egypten gezogen, wie Gott oben im 15. Capitel dem Abraham zuvor saget, sein Same soll vierhundert Jahre in Egypten wohnen, und denn allererst die Verheißung vom Lande Canaan erlangen. Nun ist’s tröstlich zu hören, daß nur siebenzig Menschen in Egypten ziehen, v. 27, und in vierhundert Jahren zu einem solchen großen Volke werden, daß sechsmal hundert tausend Mann, ohne Weib und Kinder, aus Egypten in das Land Canaan ziehen. Es ist auch ein fein Exempel für die Kinder, daß Joseph, ob er wohl ein Fürst war, seinem alten Vater nachreiset und sich auf das freundlichste gegen ihn stellet, v. 29.

\section{Votum}

Der treue Gott erfülle auch an unserm Leib und Seele seine Verheißungen, zu zeitlicher und ewiger Wohlfahrt; und helfe, daß alle Kinder ihre Eltern ehren, wie er es in seinem Wort befohlen hat, Amen.
\quellseite{117}

\chapter{Das 47. Kapitel}

\section{Vorrede}

Himmel und Erde, auch aller Kaiser, Könige und Fürsten Schlösser vermögen nicht, daß ihm Gott daraus eine Wohnung mache; aber in dem Menschen, der sein Wort hält, will Gott wohnen, sammt seinem Sohne und dem heiligen Geiste. O, ein herrlicher Trost! Die Gottes Wort glauben und darnach christlich leben, sollen der ganzen heiligen Dreifaltigkeit Tempel und Ruhhaus sein; da doch wir Menschen von Natur arme Sündenwürmlein, Kinder des Zorns, leibeigene Knechte Satanä sein und aller Teufel Fußschemel ewig sein sollten. Wer wollte nun nicht Gottes Wort gerne hören, fleißig lernen, mit Freuden glauben und christlich gebrauchen? Denn sein wir Behausungen Gottes, so wird auch Gott seine Wohnung wider allen Sturm und Anlauf schützen, wird uns in seiner Noth verlassen und wird uns im ewigen Leben wieder beherbergen, wie wir ihn beherberget haben in unsern Herzen. Und darauf höret an das 47. Kapitel.

\section[Jakob wohnet im Lande Gosen. Die Theurung ist]{Jakob wohnet im Lande Gosen. Die Theurung ist groß in Egypten}

a Wohnen) zur Herberge, Gast sein, Fremdling sein. Non cives aut domestici hujus mundi.

b Kinder) Quia nos senes mali propter pueros omnibus bonis fruimur. Wir alten Narren essen mit den Kindern, nicht sie mit uns. Ipsi Domini, nos procuratores.

\section{Summarium}

Das ist ein fein Exempel einer schönen fürstlichen Tugend am Könige Pharao, daß er um Josephs willen seinem ganzen Geschlecht hilft, V. 6. Damit beweiset Gott wiederum, obschon die gemeinen Weltstraßen über Böse und Gute gehen, daß dennoch die Frommen einen Vortheil haben sollen.

Darnach ist wohl zu merken, wie der alte Jakob von seinem und unser aller Leben auf Erden für dem Könige Pharao redet, V. 9. Daß es nur eine kurze Wallfahrt sei, da wir keine bleibende Stätte haben, und dennoch allerlei Unglück versuchen und tragen müssen. Darum sollen wir immerdar ein Sehnen und Verlangen haben nach dem rechten Vaterlande, daß wir desselben in Geduld und gutem Gewissen erwarten.
\quellseite{118}
* Neiget sich (eig) er lag im Bette krumm, richtet sich hoch auf, neiget sich zu Häupten, betet und danket Gott; dieweil that Joseph den Eid.

\section{Summarium}

Es ist eine Historia, in welcher uns eine fürstliche, herrliche Tugend wird angezeigt: obwohl ganz Egypten der Theurung halben an Geld und Vieh entblößet, die Aecker verkauft und der Fünfte von dem Gewächs dem Könige bringet wird; daß dennoch die Priester nicht beschweret, sondern denselbigen von Hofe ihr Aufenthalt gereicht wird, V. 26. Denn dies soll das fürnehmste Werck der frommen Obrigkeit sein, daß sie ob denen halten, so zu Kirchenämtern gebraucht werden. Was aber Jakob begehrt, man soll ihn begraben im Lande Canaan, V. 29, geschieht zum Anzeigen, daß er fest an der Verheißung Gottes gehalten: daß Land Canaan soll seines Samens eigen sein, ob es sich gleich lang verziehen würde.

\section{Votum}

Gott der Vater aller Barmherzigkeit verleihe, daß wir das kurze mühselige Leben recht gebrauchen und beschließen, damit wir hernach in Ewigkeit mit ihm leben in himmlischer Freude und Wonne, um JEsu Christi Verdienst willen, Amen.

\chapter{Das 48. Kapitel}

\section{Vorrede}

Die Kinder sollen ihre Eltern im Herzen aus Demuth hoch und groß, als Gottes Statthalter achten und halten, weil ihnen Gott die Ehre gibt und sie gleich neben sich an die Stelle setzt. Sie sollen ihnen von Gottes wegen schuldigen Gehorsam leisten, und sollen thun, was ihnen lieb, und lassen, was ihnen leid ist, sofern es nicht wider Gott ist; sie sollen sich mit Geduld von ihnen strafen und züchtigen lassen; sie sollen sie im Fall der Noth unterhalten und versorgen; und da sie Alters oder Krankheit halben wunderlich sind, sollen sie ihre Fehle und Mängel zu gut halten, dieselben mit Sems und Japhets Mantel bedecken. In Summa: die Kinder sollen sich nach ihren Eltern richten, und sie allwege für klüger, weiser, verständiger, denn sich halten, und sollen alle
\quellseite{119}
ihre Befehle dahin deuten, daß sie es gut mit ihnen meinen. Also soll man eben Vater und Mutter mit Worten und Geduld, auf daß ihr Segen über die Kinder komme, Sirach 3. Davon unterrichtet auch das 48. Kapitel.

\section{Jakobs Testament wegen Ephraim und Manasse}

a Stück heißt im Hebräischen Sichem, und dieselbe Stadt meinet er hie.

\section{Summarium}

Dies ist ein fein Exempel, wie Joseph seinem kranken Vater nachreiset und ihn besucht, B. 1. Und begehret seines Segens über seine Kinder, daß sie nach Jakobs, Isaaks und Abrahams Namen genennet werden, das ist, daß sie der Verheißung theilhaftig werden, welche von Gott
\quellseite{120}
den heiligen Vätern, Jakob, Isaak und Abraham, geschehen war. Daß aber Jakob den erstgeborenen Sohn an des jüngsten Statt stellt, V. 14., geschieht wie oben gemeldet, auf daß man lerne, Gott wolle mit uns nicht nach Verdienst, noch nach dem äußerlichen Ansehen, sondern nach lauter Gnade handeln.

\section{Votum}

Der treue liebe GOtt erleuchte und regiere alle Kinder, daß sie ihre Eltern ehren, damit sie GOtt wieder ehre; und segne uns alle, als unser lieber Vater, beide zu Leibes und Seelen Nothdurft und Wohlfahrt, Amen.

\chapter{Das 49. Kapitel}

\section{Vorrede}

Wie geringe Hoffnung sein kann, wenn der Same verdirbt und Schaden nimmt, daß eine gute Ernte darauf erfolgen könne: also kann auch nichts Gutes dem Vaterland und gemeinem Nutz daraus erwachsen, wenn die zarte Jugend, (welche die Alten des gemeinen Vaterlandes Hoffnung genannt) in der Blüthe verderbet und übel auferzogen wird. Derwegen hat Jakob der Patriarch seine Kinder nicht allein treulich die Zeit seines Lebens, sondern auch auf seinem Todbette unterrichtet, wie sie sich gottselig halten sollen. Dergleichen haben gethan David, und der alte Tobias an ihrem Todbette, 1 Buch der Könige 2., Tob. 4. Auch die Mutter der sieben Söhne, die Antiochus der Tyrann jämmerlich, um Gottes Worts willen, hat hinrichten lassen, 2 Maccab. 7. Wohl nun allen Eltern, die ihre Kinder aufziehen in der Zucht und Vermahnung zu dem Herrn. Dies erinnert auch das 49. Kapitel.

\section{Jakobs Weissagung und Abschied}

a Ruben sollte der Erstgeburt Würde haben, nämlich das Priesterthum und Königreich. Nun aber wird beides von ihm genommen, und Levi das Priesterthum, und Juda das Königreich gegeben. \unsicher{Dieß bedeutet die Synagoga, die das Bette Jakob, das ist, die Schrift befleckt mit falscher Lehre, darüber sie verloren hat Priesterthum und Königreich.}

b Die fängt an der Segen von Christo, der von Juda geboren sollte werden. Und heißt ihn Silo, das ist, der glückselig sein und frisch durchbringen sollte mit Geist und Glauben, daß zuvor durch Werke faul und unselig Ding war. Darum nennen wir Silo einen Held. Denn der vorige Theil dieses Segens betrifft den König David; und ist sonst in allen Segen nichts mehr von Christo, sondern alles andere ist von zeitlichem Heil, das den Kindern Israel gegeben ist, als daß Sebulon sollte am Meer wohnen, bis gen Sidon, und Isaschar mitten im Lande vom Meer wohnen und doch zinsbar gewesen ist den Königen von Assyrien.

c Den Segen Dan hat Simson erfüllet.
\quellseite{121}

\section{78. Das 1. Buch Mose 49}

d Gad hat seinen Segen ausgerichtet, da sie vor Israel herzogen.

e Asser hat gut Getreideland innen gehabt.

f Naphtali Segen ist erfüllet durch Debora und Barak.

\section{Summarium}

Hier ist eine Weissagung auf alle Stämme der Kinder Israel, wie es ihnen bis auf die Zukunft Christi gehen soll. Und ist hie sonderlich zu merken, daß Gott den Ehebruch an Ruben, B. 4., und den Todschlag an Simeon und Levi B. 6., so hart strafet, daß alle ihre Nachkommen nach ihnen solches entgelten, und die Strafe nun tragen müssen. Solches Exempel sollen alle Väter und Mütter fleißig merken, und ihrer Kinder halb sich in Gottes Furcht halten, auf daß nicht der Fluch auf sie erbe, welchen die Eltern mit ihren Sünden verdienet haben.

Die fürnehmste Weissagung aber ist von Juda, welcher Stamm zum Regiment durch David kommen, und dabei blieben ist bis auf die Maccabäer. Darum gehet die Prophezeiung also: Wenn es mit dem Königreich Juda werde ein Ende sein, und die Propheten aufhören, daß man alsdann des HErrn Christi gewarten soll. Da sollen die Heiden auch zum Reich Gottes kommen; welches Reich soll so herrlich sein, daß man Esel an die Weinstöcke binden, und die Kleider in Wein waschen soll, B. 11. Solches ist auf das einfältigste zu verstehen von der Lehre, welche im Neuen Testament geprediget wird. Da gehet Gottes Gnade in vollem Schwang, und auf das allerreichlichste, Vergebung der Sünde, und ewiges Leben wird allen Menschen angeboten, wie die Gleichnisse im Evangelio vom Hochzeit-Mahl ausweisen. Es wird aber auch an andern Orten in den Propheten das Evangelium dem Wein verglichen, dadurch die Herzen fröhlich werden, und der alte Adam in uns getödtet wird.

g Der Segen Joseph geht auf das Königreich Israel und ist ganz vom leiblichen Regiment gesagt, daß die Töchter (das ist die Städte im Lande) wohl regieret worden zeitlich, und viele Propheten und große Leute zu Ephraim hatten. Und wiewohl sie oft angefochten worden, gewannen sie doch. Und dieses Königreich war im Geschlecht Ephraim. Also bleibt der geistliche Segen und Reich auf Juda, und das leibliche Reich auf Ephraim.

h Benjamins Segen hat erfüllet der König Saul und die Bürger zu Gaba.

Kap. 50. v. 1. Da fiel Joseph auf seines Vaters Angesicht, und weinete *über ihm, und küssete ihn.

\section{Summarium}

In des Josephs zweier Söhne Segen, welche nach der Verheißung Juda von Christo das Größte haben. Denn das Königreich Israel ist mächtig und groß gewesen, und hat Herrn und
\quellseite{122}
köstliche Steine, das ist Könige und herrliche Propheten gehabt. Aber da sie denselben nicht folgen wollten, mußte ihr Land von Heiden verwüstet, und sie gefangen weggeführt werden.

Und zum Exempel, daß wir Gottes Wort und die Edelsteine seiner Kirchen, die frommen treuen Lehrer sollen lernen hoch halten, und ihnen folgen, sonst wird es nicht ungestraft bleiben.

\section{Votum}

Wir danken Gott, daß wir sein Wort und treue Kirchenlehrer haben; der wolle sie und uns segnen und behüten, und helfen, daß uns sein Wort ein Lehr-, Nähr- und Wehr-Wort sei allezeit, Amen.

\chapter{Das 50. Capitel}

\section{Vorrede}

Wie uns der gnädige und gütige Gott seine Verheißungen treulich hält, also sollen wir uns gegen ihn nach seinen Geboten gehorsamlich verhalten. Nun vergibt uns Gott aus Barmherzigkeit um Jesu Christi seines lieben Sohns willen, alle unsere manchfältige und schwere Sünden. O eine große unaussprechliche Gnade und Wohlthat! Derowegen sollen wir auch unsern Nächsten vergeben, was er uns zu Leid oder Schaden, mit Worten oder Werken gethan hat, und sollen ihm auch dienen und rathen nach unserm höchsten Vermögen; denn also lehret uns der Herr Christus, Luc. 6. und sagt: Seid barmherzig, wie auch euer Vater barmherzig ist. Und Matth. 5. spricht er: Liebet eure Feinde, thut wohl denen die euch hassen, bittet für die so euch beleidigen und verfolgen, auf daß ihr Kinder seid eures Vaters im Himmel. Davon höret auch an das 50. und letzte Kapitel.

\section{Jakob wird begraben, Joseph stirbt}
